\PassOptionsToPackage{unicode}{hyperref}
\PassOptionsToPackage{hyphens}{url}
\PassOptionsToPackage{dvipsnames,svgnames,x11names}{xcolor}
\documentclass[10pt,letterpaper]{article}

\usepackage{xcolor}
\usepackage[margin=2.0cm]{geometry}
\usepackage{lmodern}
\usepackage{amsmath,amssymb,amsfonts,mathtools}
\usepackage{textcomp}
\usepackage{microtype}
\usepackage{parskip}
\usepackage{longtable,booktabs,array,calc}
\usepackage{caption}
\captionsetup[table]{skip=6pt,font=small,labelfont=bf}
\captionsetup[figure]{font=small,labelfont=bf}
\newcounter{none}
\usepackage{etoolbox}
\usepackage{graphicx}
\usepackage{fancyvrb}
\newcommand{\VerbBar}{|}
\newcommand{\VERB}{\Verb[commandchars=\\\{\}]}
\DefineVerbatimEnvironment{Highlighting}{Verbatim}{commandchars=\\\{\},fontsize=\footnotesize}
\RecustomVerbatimEnvironment{verbatim}{Verbatim}{fontsize=\footnotesize}
\newenvironment{Shaded}{}{}
\newcommand{\AlertTok}[1]{\textcolor[rgb]{0.80,0.00,0.00}{\textbf{#1}}}
\newcommand{\AnnotationTok}[1]{\textcolor[rgb]{0.38,0.63,0.69}{\textbf{\textit{#1}}}}
\newcommand{\AttributeTok}[1]{\textcolor[rgb]{0.49,0.56,0.16}{#1}}
\newcommand{\BaseNTok}[1]{\textcolor[rgb]{0.25,0.63,0.44}{#1}}
\newcommand{\BuiltInTok}[1]{\textcolor[rgb]{0.00,0.50,0.00}{#1}}
\newcommand{\CharTok}[1]{\textcolor[rgb]{0.25,0.44,0.63}{#1}}
\newcommand{\CommentTok}[1]{\textcolor[rgb]{0.38,0.63,0.69}{\textit{#1}}}
\newcommand{\CommentVarTok}[1]{\textcolor[rgb]{0.38,0.63,0.69}{\textbf{\textit{#1}}}}
\newcommand{\ConstantTok}[1]{\textcolor[rgb]{0.53,0.00,0.00}{#1}}
\newcommand{\ControlFlowTok}[1]{\textcolor[rgb]{0.00,0.44,0.13}{\textbf{#1}}}
\newcommand{\DataTypeTok}[1]{\textcolor[rgb]{0.56,0.13,0.00}{#1}}
\newcommand{\DecValTok}[1]{\textcolor[rgb]{0.25,0.63,0.44}{#1}}
\newcommand{\DocumentationTok}[1]{\textcolor[rgb]{0.73,0.13,0.13}{\textit{#1}}}
\newcommand{\ErrorTok}[1]{\textcolor[rgb]{0.80,0.00,0.00}{\textbf{#1}}}
\newcommand{\ExtensionTok}[1]{#1}
\newcommand{\FloatTok}[1]{\textcolor[rgb]{0.25,0.63,0.44}{#1}}
\newcommand{\FunctionTok}[1]{\textcolor[rgb]{0.02,0.16,0.49}{#1}}
\newcommand{\ImportTok}[1]{\textcolor[rgb]{0.00,0.50,0.00}{\textbf{#1}}}
\newcommand{\InformationTok}[1]{\textcolor[rgb]{0.38,0.63,0.69}{\textbf{\textit{#1}}}}
\newcommand{\KeywordTok}[1]{\textcolor[rgb]{0.00,0.44,0.13}{\textbf{#1}}}
\newcommand{\NormalTok}[1]{#1}
\newcommand{\OperatorTok}[1]{\textcolor[rgb]{0.40,0.40,0.40}{#1}}
\newcommand{\OtherTok}[1]{\textcolor[rgb]{0.00,0.44,0.13}{#1}}
\newcommand{\PreprocessorTok}[1]{\textcolor[rgb]{0.74,0.48,0.00}{#1}}
\newcommand{\RegionMarkerTok}[1]{#1}
\newcommand{\SpecialCharTok}[1]{\textcolor[rgb]{0.25,0.44,0.63}{#1}}
\newcommand{\SpecialStringTok}[1]{\textcolor[rgb]{0.73,0.40,0.53}{#1}}
\newcommand{\StringTok}[1]{\textcolor[rgb]{0.25,0.44,0.63}{#1}}
\newcommand{\VariableTok}[1]{\textcolor[rgb]{0.10,0.09,0.49}{#1}}
\newcommand{\VerbatimStringTok}[1]{\textcolor[rgb]{0.25,0.44,0.63}{#1}}
\newcommand{\WarningTok}[1]{\textcolor[rgb]{0.38,0.63,0.69}{\textbf{\textit{#1}}}}

\usepackage{newunicodechar}
\newunicodechar{µ}{\ifmmode\mu\else\textmu\fi}
\newunicodechar{°}{\ifmmode^\circ\else\textdegree\fi}
\newunicodechar{·}{\ifmmode\cdot\else\textperiodcentered\fi}
\newunicodechar{²}{\ensuremath{^{2}}}
\newunicodechar{³}{\ensuremath{^{3}}}
\newunicodechar{¹}{\ensuremath{^{1}}}
\newunicodechar{⁰}{\ensuremath{^{0}}}
\newunicodechar{⁴}{\ensuremath{^{4}}}
\newunicodechar{⁵}{\ensuremath{^{5}}}
\newunicodechar{⁶}{\ensuremath{^{6}}}
\newunicodechar{⁷}{\ensuremath{^{7}}}
\newunicodechar{⁸}{\ensuremath{^{8}}}
\newunicodechar{⁹}{\ensuremath{^{9}}}
\newunicodechar{⁺}{\ensuremath{^{+}}}
\newunicodechar{⁻}{\ensuremath{^{-}}}
\newunicodechar{₀}{\ensuremath{_{0}}}
\newunicodechar{₁}{\ensuremath{_{1}}}
\newunicodechar{₂}{\ensuremath{_{2}}}
\newunicodechar{₃}{\ensuremath{_{3}}}
\newunicodechar{₄}{\ensuremath{_{4}}}
\newunicodechar{ₙ}{\ensuremath{_{n}}}
\newunicodechar{ν}{\ensuremath{\nu}}
\newunicodechar{Δ}{\ensuremath{\Delta}}
\newunicodechar{Ω}{\ensuremath{\Omega}}
\newunicodechar{α}{\ensuremath{\alpha}}
\newunicodechar{β}{\ensuremath{\beta}}
\newunicodechar{δ}{\ensuremath{\delta}}
\newunicodechar{ε}{\ensuremath{\varepsilon}}
\newunicodechar{θ}{\ensuremath{\theta}}
\newunicodechar{κ}{\ensuremath{\kappa}}
\newunicodechar{λ}{\ensuremath{\lambda}}
\newunicodechar{π}{\ensuremath{\pi}}
\newunicodechar{ρ}{\ensuremath{\rho}}
\newunicodechar{σ}{\ensuremath{\sigma}}
\newunicodechar{τ}{\ensuremath{\tau}}
\newunicodechar{→}{\ensuremath{\rightarrow}}
\newunicodechar{∎}{\ensuremath{\blacksquare}}
\newunicodechar{−}{\ensuremath{-}}
\newunicodechar{∧}{\ensuremath{\wedge}}
\newunicodechar{∫}{\ensuremath{\int}}
\newunicodechar{≈}{\ensuremath{\approx}}
\newunicodechar{≤}{\ensuremath{\leq}}
\newunicodechar{≥}{\ensuremath{\geq}}
\newunicodechar{⊆}{\ensuremath{\subseteq}}
\newunicodechar{∞}{\ensuremath{\infty}}
\newunicodechar{š}{\ifmmode\check{s}\else\v{s}\fi}
\newunicodechar{✓}{\ensuremath{\checkmark}}
\newunicodechar{≠}{\ensuremath{\neq}}
\newunicodechar{↔}{\ensuremath{\leftrightarrow}}
\newunicodechar{≪}{\ensuremath{\ll}}
\newunicodechar{↛}{\ensuremath{\nrightarrow}}
\newunicodechar{⇒}{\ensuremath{\Rightarrow}}
\newunicodechar{∪}{\ensuremath{\cup}}
\newunicodechar{⟨}{\ensuremath{\langle}}
\newunicodechar{⟩}{\ensuremath{\rangle}}
\newunicodechar{・}{\textperiodcentered}

\usepackage{xurl}
\usepackage{bookmark}
\usepackage{hyperref}

\setcounter{secnumdepth}{-\maxdimen}
\providecommand{\tightlist}{\setlength{\itemsep}{1pt}\setlength{\parskip}{0pt}}

\definecolor{openieAmber}{HTML}{D97706}
\definecolor{openieAmberBg}{HTML}{FFFBEB}
\definecolor{openieAmberBorder}{HTML}{F59E0B}
\definecolor{openieSlate}{HTML}{334155}
\definecolor{openieDark}{HTML}{0F172A}
\definecolor{openieCodeBg}{HTML}{F8FAFC}
\definecolor{openieCodeBorder}{HTML}{CBD5E1}
\definecolor{openieBlue}{HTML}{0057B8}

\hypersetup{
  pdftitle={RCE Compile: A Law of Compilation, Tested as a Normal Form Against Real Compilers},
  pdfauthor={David Charlot, Open Interface Engineering},
  colorlinks=true,
  linkcolor=openieBlue,
  citecolor=openieBlue,
  urlcolor=openieBlue,
  pdfcreator={XeTeX / Tectonic — Open Interface Engineering}
}

\urlstyle{same}
\setlength{\emergencystretch}{3em}
\hfuzz=6pt
\sloppy

\newcommand{\openiefig}[3]{%
  \IfFileExists{#1}{%
    \includegraphics[width=#2\linewidth,height=0.70\textheight,keepaspectratio]{#1}%
  }{%
    \fcolorbox{openieCodeBorder}{openieCodeBg}{%
      \parbox{0.92\linewidth}{\footnotesize\ttfamily Figure asset embedded in archival PDF; interactive/vector source: \url{#3}}%
    }%
  }%
}

\title{\vspace{-1.2em}\textbf{RCE Compile: A Law of Compilation, Tested as a Normal Form Against Real Compilers}}
\author{\textbf{David Charlot, Open Interface Engineering}\\[0.15em]
{\small Research study · draft $\cdot$ Open Interface Engineering $\cdot$ \href{https://research.openie.dev}{research.openie.dev}}}
\date{\small Revision \texttt{395f081} (4 November 2026) $\cdot$ Web: \href{https://research.openie.dev/papers/rce-compile/}{research.openie.dev/papers/rce-compile/} $\cdot$ Archival LaTeX: \href{https://research.openie.dev/pdfs/latex/rce-compile.pdf}{/pdfs/latex/rce-compile.pdf} (\href{https://research.openie.dev/pdfs/latex/rce-compile.tex}{.tex})}

\begin{document}
\maketitle
\vspace{-0.6em}

% OpenIE Measurement Honesty Callout Box
\noindent\fcolorbox{openieAmberBorder}{openieAmberBg}{%
  \parbox{\dimexpr\linewidth-2\fboxsep-2\fboxrule\relax}{%
    \small\color{openieDark}%
    \textbf{Measurement.} Research study, draft, not a journal final. Measured results in Section~8 cover byte comparisons of compiler outputs across three hosts (\texttt{E1}), erasure and description-length counts alongside Regime-B Linux \texttt{x86-64} \texttt{RAPL} core/package \textbf{\texttt{reported\_j}} compilation cost and optimization payback horizons across \texttt{20} exact-integer kernels and \texttt{14} compiler tiers (\texttt{E4} and \texttt{E4-Energy}: Python$\to$\texttt{clang -O3} median payback \texttt{2,490} runs; \texttt{-O0}$\to$\texttt{-O3} \texttt{53,464--102,720} runs), and De~Bruijn positional name-free hashing across \texttt{245} programs and \texttt{500} verified semantic mutants (\texttt{E5} and \texttt{E5-Scale}: \texttt{0 / 500} false hits; \texttt{167.8$\times$} \texttt{clang -O3} and \texttt{448.5$\times$} \texttt{rustc -O3} rebuild \texttt{reported\_j} reduction), plus the 12-compiler normal-form audit (\texttt{R..G}). No hardware rail or wall wattmeter was attached (\texttt{measured\_j = None}, \texttt{board\_metered = false}); every external source states how it was verified.
  }%
}

\vspace{6pt}
% OpenIE Amber Left-Rule Summary Quote
\noindent{\color{openieAmber}\vrule width 3pt}\hspace{8pt}%
\parbox{\dimexpr\linewidth-11pt\relax}{%
  \small\itshape\color{openieSlate}%
  RCE Compile (a research project in the Reality Code Engine repository) asks whether one statement of what compilation is holds for every compiler, independent of the Reality Code Engine itself. Checked against twelve compiler families, the candidate holds as a standard a compiler can be measured against, not as a description of what compilers do. This study states that standard, the evidence for and against each of its seven parts, a measurable replacement for its convergence function, seven hypotheses with their falsifiers, and the experiments run so far: byte comparisons of compiler outputs on three machines, erasure counts alongside Linux Running Average Power Limit (RAPL) reported\_j compilation cost and optimization payback horizons across 20 kernels and 14 compiler tiers, and name-free De Bruijn reuse hashing tested across 245 programs and 500 verified semantic mutants.
}
\vspace{6pt}

{
\small
\setcounter{tocdepth}{2}
\tableofcontents
}
\vspace{4pt}
\hrule
\vspace{8pt}

\emph{Revision
\href{https://github.com/openIE-dev/research-openie-web/commits/main/src/content/papers/rce-compile.md}{\texttt{395f081}},
4 November 2026. Research study, draft. Project repository:
\href{https://github.com/openIE-dev/rce/tree/main/compile}{openIE-dev/rce,
\texttt{compile/}}, which holds the charter, the hypotheses and the
experiments this study reports. Interactive observatory:
\href{https://research.openie.dev/living/rce-compile/}{\texttt{/\allowbreak{}living/\allowbreak{}rce-\allowbreak{}compile/\allowbreak{}}}.}

\section{Abstract}\label{abstract}

RCE Compile is a research project that asks whether a single statement
of what compilation is, here called a Law of Compilation, holds for any
compiler, independent of the Reality Code Engine (RCE) in whose
repository the project lives. The candidate statement says that
compilation is meaning-preserving lowering between canonical
representations, ranked by a cost whose floor is Landauer's limit and
whose target is the shortest equivalent program, gated by verification
and compounded by reuse. This study checks that candidate against twelve
families of production compilers and the research literature behind
them. The candidate fails as a description: no compiler surveyed hashes
a name-free representation, lowering in shipping compilers is neither
total nor shown to preserve meaning, and no compiler accounts for
energy. It holds up as a normative framework once three clauses are
rewritten, and the study states it in that form: a normal form with two
bounds, plus a measured distance from them. The convergence function
cannot be computed or measured as stated, and the study replaces it with
two counts that can be computed and compared with measured energy. Seven
hypotheses with falsifiers follow, with the experiments run so far and
the protocols of two more. Compiler outputs were compared byte for byte
on three machines; some differences survive the toolchains' path
remapping and have no traced cause, so the hypothesis that lowering is
deterministic is neither supported nor falsified. Erasure and
description length were first counted for 96 builds on a machine with no
energy counter, and a pre-registered Regime-B extension then read Linux
Running Average Power Limit (RAPL) on-chip \texttt{reported\_\allowbreak{}j}
(\texttt{board\_\allowbreak{}metered\ =\ false}, no wall wattmeter) across 20
deterministic integer and fixed-point kernels and 14 compiler and
runtime tiers (280 of 280 builds checksum-verified): lowering from
Python to \texttt{clang\ -\allowbreak{}O3} recovers its compilation joules after a
median of 2,490 runs, \texttt{-O0} to \texttt{-O3} pays back after a
median of 53,464 to 102,720 runs, \texttt{-O2} to \texttt{-O3} never
recovers its extra compile energy on 4 to 6 of the 20 kernels, and time
multiplied by power explains far more Regime-B variance than static
intermediate-representation size (adjusted \(R^2\) 0.9999 against
0.1154). A name-free hash, the part of the standard that no compiler
surveyed has, was built first over binary objects and WebAssembly
modules (16 kernels and 4 programs) and then scaled across 185 RCE
Forms, 60 C, Rust and Python kernels, and 500 test-vector-verified
semantic mutants: positional De Bruijn hashing produced 0 of 500
(0.00\%) false-positive collisions while cutting cold
\texttt{clang\ -\allowbreak{}O3} and \texttt{rustc\ -\allowbreak{}O3} compilation
\texttt{reported\_\allowbreak{}j} by 167.8× and 448.5× on renamed hits, whereas naive
identifier-stripping collided on 150 of 150 (100.0\%) variable-binding
swaps. Almost every external source was seen only as a search result,
and each citation says so.

\textbf{In plain words}

\begin{enumerate}
\def\labelenumi{\arabic{enumi}.}
\tightlist
\item
  The question is whether all compilers obey one rule. Real compilers do
  not, so the rule is restated as a standard that any compiler can be
  measured against.
\item
  The standard has seven parts. Most compilers have some of them. None
  keeps track of energy, and none proves that it translates every
  program correctly.
\item
  Heavy compiler optimization spends chip energy up front to save chip
  energy per run. Across 20 test kernels on Linux chip counters
  (\texttt{reported\_\allowbreak{}j}, not a wall plug meter), compiling to
  \texttt{clang\ -\allowbreak{}O3} instead of running Python pays back after about
  2,500 runs, while stepping from \texttt{-O2} to \texttt{-O3} never
  pays back on up to 30\% of the kernels.
\item
  A fingerprint of a program that replaces variable names by their
  binding positions --- rather than merely erasing names --- recognised
  renamed copies across 245 programs, avoided 168× to 449× of rebuild
  chip joules on warm hits, and made zero false matches across 500
  bug-injected mutants where naive name-erasing collided on every
  swapped variable pair.
\end{enumerate}

\section{Terms and evidence marks}\label{terms-and-evidence-marks}

Every acronym on this page is either expanded where it is first used or
listed here.

\begin{itemize}
\tightlist
\item
  AI, API: artificial intelligence; application programming interface
\item
  AIR, ZIR, MIR: intermediate representations of the Zig compiler
  (analysed, untyped and machine-level)
\item
  ANF: A-normal form, a program form in which every intermediate value
  is named
\item
  APU, ARX, BST: accelerated processing unit; add-rotate-xor cipher
  structure; binary search tree
\item
  ASR: Abstract Semantic Representation, the intermediate representation
  of LFortran
\item
  AST: abstract syntax tree
\item
  BF16, FP16, FP8, FP4: 16-, 8- and 4-bit floating-point formats
\item
  CI, CRC, CRC32: continuous integration, the checks run on every change
  to a repository; cyclic redundancy check (32-bit IEEE 802.3
  polynomial)
\item
  CLIF: the intermediate representation of the Cranelift code generator
\item
  CMOS: complementary metal-oxide-semiconductor, the usual chip
  technology
\item
  CN101, MAX, XY.AI, MMA-Sim: product, platform, company and tool names
  used as they are
\item
  CNF: conjunctive normal form, a standard form for logical formulas
\item
  COBOL, SPARK, GHC, GCC, LLVM, MLIR, XLA, TVM, CUDA: names of languages
  and compilers (COBOL; SPARK, a verifiable subset of Ada; the Glasgow
  Haskell Compiler; the GNU Compiler Collection; the LLVM compiler
  infrastructure and its multi-level representation MLIR; the XLA and
  TVM compilers for tensor programs; NVIDIA's GPU platform)
\item
  CPU, GPU, FPGA: central processing unit, graphics processing unit,
  field-programmable gate array
\item
  DDC: diverse double-compiling, a test that rebuilds a compiler with an
  independent compiler and compares the bits
\item
  DRAM, SRAM, DP: dynamic and static random-access memory; dynamic
  programming
\item
  DWARF: the standard debug-information format
\item
  ELF: Executable and Linkable Format, the binary format of Linux
  executables and object files
\item
  EMI: equivalence modulo inputs, a compiler-testing method that changes
  code a given input never runs and checks the output stays the same
\item
  FNV, IQR, JS, JSON: Fowler--Noll--Vo hash; interquartile range;
  JavaScript; JavaScript Object Notation
\item
  GEMM: general matrix multiplication; cuBLAS and cuDNN are NVIDIA's
  libraries for it and for neural-network kernels
\item
  GIMPLE, FIR, HLFIR, SOACS: the intermediate representations of GCC, of
  LLVM Flang (two levels) and of Futhark
\item
  GNU, GNAT, GNATprove: the GNU project; the Ada compiler in GCC; the
  SPARK proof tool
\item
  GODEBUG: the environment variable that turns on debugging settings of
  the Go toolchain and runtime
\item
  GOPATH: Go's older workspace mode, used before Go modules
\item
  HAL, CNKI, ACM DL, J-STAGE: national and publisher research databases
  (France, China, the ACM, Japan)
\item
  HIR, THIR: the high-level and typed high-level representations of the
  Rust compiler, rustc
\item
  HOL4, Rocq, Lean, Coq: proof assistants (Coq was renamed Rocq)
\item
  HPC: high-performance computing
\item
  IEEE 754: the floating-point standard of the Institute of Electrical
  and Electronics Engineers; NaN: not-a-number
\item
  IR: intermediate representation, the form a compiler holds a program
  in between source and machine code
\item
  ISLE: the rule language in which Cranelift's lowering rules are
  written
\item
  LAPACK: the Linear Algebra Package, a Fortran library
\item
  LCNF: Lean Compiler Normal Form
\item
  LTO, LZ77: link-time optimisation; Lempel--Ziv 1977 sliding-window
  compression
\item
  KB, MB, GB: kilobytes, megabytes, gigabytes
\item
  MDL, MSR, NTT, OLS, OS: minimum description length, a two-part code
  length relative to a named class of models; model-specific register;
  number-theoretic transform; ordinary least squares; operating system
\item
  ML compilers: compilers for neural-network tensor programs, such as
  XLA, TorchInductor, TVM, TensorRT, Triton and IREE
\item
  MLGO, TENSAT, STOKE: a policy-driven LLVM optimisation framework, an
  equality-saturation optimiser for tensor graphs, and a stochastic
  superoptimiser
\item
  NOT FOUND: an evidence mark, defined below
\item
  NVML, RAPL: NVIDIA's management library and Intel's running average
  power limit interface, both of which report energy counters
\item
  OSS: open-source software
\item
  OTP: Open Telecom Platform, the name of the Erlang distribution and
  its release numbers
\item
  POPL, PLDI, ICFP, OOPSLA, MLSys, ISSCC, ICML, PriSC: conferences and
  workshops in which cited papers appeared
\item
  PP: the class of problems a polynomial-time randomised machine decides
  correctly more than half the time; PP-hard problems are believed to be
  far harder than NP-complete ones
\item
  PTX, PYC, RLE: NVIDIA's virtual GPU instruction set; Python compiled
  bytecode; run-length encoding
\item
  RCE: the Reality Code Engine
\item
  SHA-256, SIMD, SWAR, UTF: a cryptographic hash function; single
  instruction, multiple data; SIMD within a register; Unicode
  Transformation Format
\item
  SMT: satisfiability modulo theories; SMT solvers such as Z3 decide
  logical formulas over arithmetic, bit-vectors and memory
\item
  SSA: static single assignment form
\item
  STG, Cmm: the intermediate representations of GHC after Core
\item
  TIRx, JACC, NVIDIA, AMD, ONNX, BEAM, DIKU: project, company and system
  names used as they are
\item
  TZ, UTC, LC\_ALL: environment variables for the time zone and locale;
  UTC is Coordinated Universal Time
\item
  UB: undefined behaviour, source behaviour to which the language
  standard gives no meaning
\item
  UBC, MIT, ORNL, NERSC, LBNL, HPE: the University of British Columbia,
  the Massachusetts Institute of Technology, Oak Ridge National
  Laboratory, the National Energy Research Scientific Computing Center,
  Lawrence Berkeley National Laboratory and Hewlett Packard Enterprise
\item
  ULP: unit in the last place
\item
  UNCOL: Universal Computer Oriented Language, a 1958 proposal for one
  intermediate language for all compilers
\item
  WARP: the function-identification format of Binary Ninja, a
  reverse-engineering tool, used as a name
\end{itemize}

\textbf{Evidence marks.} Each external statement carries the mark of how
it was checked, and each entry in the reference list repeats it.
\textbf{{[}F{]}} means read in the source itself. \textbf{{[}S{]}} means
seen only in a search-result snippet or a secondary summary.
\textbf{{[}K{]}} means stated from memory and not re-checked; the
canonical location is given but was not opened. \textbf{{[}M{]}} means
measured by the project on its own machine. \textbf{{[}NOT FOUND{]}}
means searched for with no source located. Only two external items
reached {[}F{]}: an index page and a set of repository commit counts.
Most hosts that hold the primary texts (among them arxiv.org,
compcert.org, gcc.gnu.org, ziglang.org, zenodo.org, pytorch.org and the
UBC library) could not be reached from the research sessions, so an
{[}S{]} item below means ``a source exists and its abstract or snippet
says this'', not ``the paper was read''.

\section{1. Problem statement}\label{problem-statement}

A compiler takes a program in one language and produces a program in
another, usually closer to a machine. There are thousands of compilers,
built by communities that rarely share vocabulary: systems languages,
scientific computing, legacy business languages, high-integrity
languages, functional languages and the compilers for neural-network
tensor programs. Each community has its own account of what a correct
and good compiler is. RCE Compile asks whether one statement covers all
of them, independent of RCE, and if not, what the nearest statement is
that does.

The candidate under test came out of an earlier discussion with an
assistant system (Grok). It says: \textbf{compilation is
semantics-preserving lowering between canonical representations, ranked
by a cost whose floor is Landauer's limit and whose target is
Kolmogorov-minimal, gated by verification, and compounded by reuse.}
Three questions follow. Does it describe the compilers that exist? If
not, can it be restated so that it is true of a correct compiler and
checkable against a real one? And can its physical and mathematical
terms be measured at all? The discussion that proposed the candidate
also made 55 factual claims about the field; those are checked in
\href{https://research.openie.dev/papers/rce-compile/\#appendix-a-claims-checked}{Appendix
A}.

\section{2. The candidate and its
refinement}\label{the-candidate-and-its-refinement}

\subsection{2.1 Why the candidate fails as a
description}\label{why-the-candidate-fails-as-a-description}

The candidate makes three kinds of claim: about representation
(``canonical''), about lowering (``semantics-preserving''), and about
cost (``floor is Landauer, target is Kolmogorov-minimal''). Each fails
as a description of shipping compilers, for reasons given in full in §3
and §4:

\begin{itemize}
\tightlist
\item
  \textbf{No compiler surveyed hashes a name-free intermediate
  representation.} Every one keys its identity and its caches on named
  or textual content.
\item
  \textbf{Lowering is neither total nor proved to preserve meaning.}
  Fuzzers have found more than 325 bugs (Csmith) and 1,602 bugs (the EMI
  project) in production C compilers {[}1{]} {[}2{]}, and Zig's default
  self-hosted backend fails 24 of its 2,008 behaviour tests {[}3{]}.
\item
  \textbf{No compiler accounts for energy.} Every cost model found is in
  cycles, bytes or node counts.
\item
  \textbf{The convergence function cannot be computed or measured as
  stated.} Kolmogorov complexity is uncomputable; even the minimum heat
  of one fixed circuit is PP-hard to compute {[}4{]}; and today's
  hardware sits \(10^6\) to \(10^{11}\) times above the Landauer floor.
\end{itemize}

A statement that every compiler is claimed to obey, and that real
compilers measurably break, is falsified as a descriptive law. No
peer-reviewed source in English, Chinese, Russian or Japanese was found
that names or proves a single law of compilation {[}NOT FOUND{]}; an
English search for the phrase returns only the copyright sense of
``compilation'' {[}5{]}.

\subsection{2.2 The refinement: a normal form with two bounds and a
measured
distance}\label{the-refinement-a-normal-form-with-two-bounds-and-a-measured-distance}

The candidate holds up much better as a normative framework: a statement
of what a correct, well-ranked compiler must satisfy, against which
every real compiler can be placed and its shortfall measured. Three
clauses change. ``Semantics-preserving'' becomes a \textbf{declared
contract}, because the established notion is refinement under a
precondition and the precondition differs by compiler. ``Canonical''
becomes an \textbf{operation} the framework applies to any
representation, because no compiler supplies it. And ``Landauer floor
and Kolmogorov target'' becomes \textbf{two computable counts plus a
measurement}. The result is a normal form in the sense of Hoare, He and
Sampaio (§6), with two bounds and a measured distance from them, not a
law every compiler already obeys.

\subsection{2.3 The normative statement}\label{the-normative-statement}

The project's charter states the standard as follows; the text is quoted
from it unchanged {[}6{]}:

\begin{quote}
\textbf{Compiler Form.} A compiler is a partial function from declared
inputs (source, flags, toolchain, and any build path after remapping) to
a chain of representations R₀ → \ldots{} → Rₙ. Each step either stops or
produces an output that a named checker K accepts as preserving a
declared contract: a property class, an equivalence level, a linking
assumption and a trusted computing base. Among accepted outputs, a
stated cost C ranks the candidates. Reuse U is keyed by a name-free hash
of the representation together with its contract. Nothing is released
unless G records K's acceptance. For every released output, the
framework reports its distance from two bounds, the bits irreversibly
overwritten per run (scaled by k\textsubscript{B}·T·ln 2) and a stated
description-length proxy, next to the energy measured in the regime the
hardware belongs to.
\end{quote}

A \emph{stop} is a lowering step that ends with a diagnostic and no
output. A \emph{declared input} is an input the compiler's contract
lists. The statement drops ``holds for any compiler'' in favour of
``every compiler can be placed against it''. The project claims a law
only for what its registered hypotheses (§7) survive. §3 sets out the
parts the statement names, and §4 the two bounds.

\section{3. The Compiler Form, R to G}\label{the-compiler-form-r-to-g}

\subsection{3.1 The seven parts}\label{the-seven-parts}

The project calls the framework the \textbf{Compiler Form}. It has seven
parts:

{\footnotesize\setlength{\tabcolsep}{3.8pt}\renewcommand{\arraystretch}{1.15}
\begin{longtable}[]{@{}
  >{\raggedright\arraybackslash}p{(\linewidth - 4\tabcolsep) * \real{0.1146}}
  >{\raggedright\arraybackslash}p{(\linewidth - 4\tabcolsep) * \real{0.1875}}
  >{\raggedright\arraybackslash}p{(\linewidth - 4\tabcolsep) * \real{0.6979}}@{}}
\toprule\noalign{}
\begin{minipage}[b]{\linewidth}\raggedright
\textbf{Part}
\end{minipage} & \begin{minipage}[b]{\linewidth}\raggedright
\textbf{Name}
\end{minipage} & \begin{minipage}[b]{\linewidth}\raggedright
\textbf{What the part requires}
\end{minipage} \\
\midrule\noalign{}
\endhead
\bottomrule\noalign{}
\endlastfoot
R & Representation & A canonical IR with a stated equivalence relation
and a name-free hash \\
L & Lowering & Each step either stops (ends with a diagnostic and no
output) or produces an output that preserves a declared contract \\
C & Cost & A stated cost that ranks equivalent candidates \\
K & Checker & A named checker that certifies the contract \\
U & Reuse & Earlier accepted results are reused, keyed by the name-free
hash and the contract \\
E & Energy & Energy is accounted for, by counting or by meter, for each
released output \\
G & Gate & Nothing is released unless the gate records the checker's
acceptance \\
\end{longtable}
}

Two refinements, from the correctness and ML-compiler literature, make
L, K and G precise enough to test.

\subsection{3.2 Four equivalence levels}\label{four-equivalence-levels}

ML compilers define correctness as closeness. XLA, TorchInductor, TVM,
TensorRT, Triton and IREE all test against a reference within an
absolute and relative tolerance (PyTorch's float32 defaults are a
relative tolerance of \(1.3\times10^{-6}\) and an absolute tolerance of
\(10^{-5}\)) {[}K{]}, because fusion, tiling, split reductions and
tensor-core accumulation reorder floating-point sums, and lowering to
FP8 or FP4 changes outputs on purpose {[}K{]}. FlexAttention's
FlashAttention-4 backward pass is reported to be nondeterministic when
block sparsity is on {[}7{]}. If ``preserving'' meant identical output
bits, these compilers would all violate L. A September 2026 preprint
argues that bit-level results depend on the accumulation structure the
compiler chooses, reconstructs that structure for cuBLAS across three
GPU generations, and gives static checkers that decide bitwise
equivalence across tens of thousands of autotuner configurations
{[}8{]}; MMA-Sim gives a bit-accurate model of tensor and matrix cores
on ten architectures {[}9{]}. Halide already shows the pattern at small
scale: a schedule cannot change results except through named
reassociating reductions {[}K{]}.

The Form therefore requires every lowering to declare one of four
\textbf{equivalence levels}, and its checker to certify that level:

{\footnotesize\setlength{\tabcolsep}{3.8pt}\renewcommand{\arraystretch}{1.15}
\begin{longtable}[]{@{}
  >{\raggedright\arraybackslash}p{(\linewidth - 4\tabcolsep) * \real{0.2241}}
  >{\raggedright\arraybackslash}p{(\linewidth - 4\tabcolsep) * \real{0.3793}}
  >{\raggedright\arraybackslash}p{(\linewidth - 4\tabcolsep) * \real{0.3966}}@{}}
\toprule\noalign{}
\begin{minipage}[b]{\linewidth}\raggedright
\textbf{Level}
\end{minipage} & \begin{minipage}[b]{\linewidth}\raggedright
\textbf{Meaning}
\end{minipage} & \begin{minipage}[b]{\linewidth}\raggedright
\textbf{What K must certify}
\end{minipage} \\
\midrule\noalign{}
\endhead
\bottomrule\noalign{}
\endlastfoot
Exact & Integer, Boolean or fixed-point; identical output bits & Bitwise
equality \\
Bitwise within an accumulation class & Floating point, with the
reduction order recorded in R; L and C choose only within the class &
Same class, then bitwise equality \\
Within a declared error bound & Floating point, with an absolute,
relative or ULP bound & The bound, by static error analysis or interval
methods, not by sample tests \\
Declared lossy & Quantisation or precision lowering & Nothing about
equality: it is a change of program, recorded and admitted by G as a
separate decision \\
\end{longtable}
}

The default is ``exact unless a checked contract says otherwise''. That
is stricter than every production ML compiler and consistent with the
bitwise-class research.

\subsection{3.3 What a gate must name}\label{what-a-gate-must-name}

The classical correctness theorem is not ``the output means the same as
the input''. CompCert's main theorem is \textbf{behaviour refinement
under a precondition}: if compilation succeeds and the source has no
undefined behaviour, every observable behaviour of the compiled code is
a behaviour of the source {[}K{]}. Patterson and Ahmed show that ``the
compiler is correct'' is ambiguous until the theorem says what target
code the output may be linked with {[}10{]}. Secure compilation adds
that a compiler can preserve meaning in CompCert's sense and still let
an attacker at the target level break abstractions of the source:
\emph{full abstraction} (source-equivalent components compile to
target-equivalent ones) and the \emph{robust property preservation}
criteria are distinct, stronger guarantees {[}11{]}. And CompCert's
proof leaves the preprocessor, assembler and linker outside it; Csmith
found silent miscompilations exactly there and none in the verified
middle end {[}1{]}.

From these, every gate in the Form must state three things:

\begin{enumerate}
\def\labelenumi{\arabic{enumi}.}
\tightlist
\item
  \textbf{The property class it preserves}: trace properties,
  hyperproperties, a robust variant, or an equivalence.
\item
  \textbf{The linking assumption}: which other code the output may be
  combined with.
\item
  \textbf{The trusted base}: the semantics, proof kernel, assembler,
  hardware models or diverse compiler the guarantee relies on.
\end{enumerate}

Fully abstract compilation usually needs run-time checks or isolation,
so a stronger property class raises cost {[}K{]}.

On 2 November 2026, before any gate was measured, the project's charter
added two statements to these three {[}6{]}:

\begin{enumerate}
\def\labelenumi{\arabic{enumi}.}
\setcounter{enumi}{3}
\tightlist
\item
  \textbf{The authority of its evidence}: the level that the evidence
  behind the checker's acceptance must have.
\item
  \textbf{The digests its verdict is bound to}: the SHA-256 (the
  \emph{digest}) of the declared inputs, of the output it releases and
  of the checker that accepted it, and the revision of the contract it
  checked. A verdict bound to a digest applies only to the exact bytes
  it was made on.
\end{enumerate}

The charter takes the levels of authority from the reconstruction
ledgers of REA (Reverse Engineer Anything), an open-source
reverse-engineering project {[}6{]}. Two levels are \emph{observed}
evidence: \emph{pinned}, in which the compiler, named by the SHA-256 of
its toolchain, is run on the declared inputs and its output compared
byte for byte; and \emph{replay}, in which an output recorded earlier
and bound by its SHA-256 is compared byte for byte. \emph{Proved}
evidence is a checker's certificate for every input of its fragment: a
proof that a proof kernel checks, a translation validator's verdict, or
an SMT solver's finding that the negated contract cannot be satisfied.
\emph{Derived} evidence follows from a stated rule that cites its
source, such as a documented gate behaviour quoted with its passage.
\emph{Inferred} evidence is the statement of a model, a heuristic or a
stochastic proposer, or a summary that cites no record of a higher
level. A gate accepts an output only when evidence of the authority it
states holds and is bound to every digest it states; a verdict bound to
other bytes, or evidence that cannot be read, releases nothing. Inferred
evidence never makes a gate accept. Observed evidence shows that a
contract held, or failed, on the inputs that were run; only proved
evidence certifies it for every input of a fragment. A gate that does
not state all five is recorded as a gate with an unstated guarantee.

\subsection{3.4 Evidence per part across
compilers}\label{evidence-per-part-across-compilers}

The survey covered Zig, Rust, Go, GCC and Clang/LLVM, Swift, Cranelift
and WebAssembly (Wasm), Julia, the Fortran compilers (gfortran, ifx,
nvfortran, LLVM Flang, LFortran), COBOL (gcobol and GnuCOBOL), Ada and
SPARK, GHC and the other functional compilers (OCaml, Lean 4, Idris 2,
Koka, Erlang, Futhark, Halide), and the ML compilers. Figure 1 and the
table below summarise the evidence part by part. Almost every cell rests
on {[}K{]} or {[}S{]}.


\begin{center}
\openiefig{fig/rce-compile-fig01-form-evidence.png}{0.94}{https://research.openie.dev/figures/rce-compile/form-evidence.svg}
\end{center}


\emph{Figure 1. Where each part of the Compiler Form is found, by
compiler family. Filled: present as the Form requires. Half: present in
part or only by configuration. Empty: absent. Erlang's code validator, a
real gate, sits inside the functional row. Qualitative summary of the
survey in §3.4 and §3.5; almost every cell rests on {[}K{]} or {[}S{]}
evidence.}

{\footnotesize\setlength{\tabcolsep}{3.8pt}\renewcommand{\arraystretch}{1.15}
\begin{longtable}[]{@{}
  >{\raggedright\arraybackslash}p{(\linewidth - 6\tabcolsep) * \real{0.0530}}
  >{\raggedright\arraybackslash}p{(\linewidth - 6\tabcolsep) * \real{0.3485}}
  >{\raggedright\arraybackslash}p{(\linewidth - 6\tabcolsep) * \real{0.3636}}
  >{\raggedright\arraybackslash}p{(\linewidth - 6\tabcolsep) * \real{0.2348}}@{}}
\toprule\noalign{}
\begin{minipage}[b]{\linewidth}\raggedright
\textbf{Part}
\end{minipage} & \begin{minipage}[b]{\linewidth}\raggedright
\textbf{Supported by}
\end{minipage} & \begin{minipage}[b]{\linewidth}\raggedright
\textbf{Qualified by}
\end{minipage} & \begin{minipage}[b]{\linewidth}\raggedright
\textbf{Falsified by}
\end{minipage} \\
\midrule\noalign{}
\endhead
\bottomrule\noalign{}
\endlastfoot
R & E-classes; Wasm's name-free binary; exact text round-trip in LLVM IR
and CLIF; normal forms (GHC Core, LFortran ASR, Lean LCNF) & Every
compiler is a chain of 3 to 6 IRs; round-trip holds only between an IR
and its own bytes & No compiler hashes a name-free IR \\
L & Morris's commuting square; CompCert; CakeML; Halide's schedule
separation & Refinement, not equality; only for UB-free sources; indexed
by linking context; ML compilers use tolerance & Csmith, EMI and Alive2
bug counts; Zig's failing behaviour tests \\
C & Cranelift e-graph extraction; equality saturation; autotuners &
Mostly local heuristics in cycles or bytes & No cost in joules; shortest
is not cheapest \\
K & Type and borrow checkers; IR verifiers; Alive2; ISLE verification;
the Lean kernel; SPARK & Checks the source more than the lowering; often
debug-only; bounded; ML checks are sample tests & No production compiler
checks all of its lowering \\
U & Content-addressed caches in Go, rustc, LLVM, Wasmtime, JAX, Triton,
TVM; CakeML bootstrapping; DDC & Keys include names and paths;
TensorRT-LLM dropped its stored engines & Name-free reuse is absent \\
E & Only research (Eel, MIT 2016) and external meters & Energy roughly
tracks run time on CPUs & Absent in every compiler surveyed \\
G & rustc, Go, Wasm validation, Erlang's validator, SPARK when a project
requires it & tsc emits despite type errors by default; LLVM and GCC
gate on post-commit testing; ML gates are fallbacks & No universal
gate \\
\end{longtable}
}

\subsection{3.5 Part-by-part findings}\label{part-by-part-findings}

\textbf{R is absent everywhere as the Form states it.} Every production
compiler lowers through a chain of IRs rather than one canonical form.
Zig goes AST → ZIR → AIR → per-backend MIR {[}3{]}; rustc goes AST → HIR
→ THIR → MIR → one of LLVM IR, Cranelift's CLIF or GCC's GIMPLE
{[}12{]}; Go goes syntax → IR → SSA → Plan 9-style instructions
{[}13{]}; GHC goes Core → STG → Cmm {[}K{]}. Canonical \emph{forms}
exist (GHC Core, LFortran's ASR, Futhark's SOACS, Lean's LCNF, Flang's
FIR and HLFIR), and the e-classes of an e-graph are canonical
\emph{sets} of equivalent terms. The closest result to the ``canonical
representations'' clause is Bryant's thesis, which treats compilation as
normalisation to a normal form {[}14{]}. Two qualifications follow.
Morris's correctness square needs only \emph{some} semantics on each
side, not canonicity {[}K{]}. And round-trip holds only between an IR
and its own serialised form (LLVM's text and bitcode, CLIF text, MLIR's
generic form), never back to source, because every compiler discards
comments, names and syntactic sugar {[}K{]}. The falsifier is plain:
GHC's interface fingerprints, OCaml's interface digests, rustc's stable
hashes, Go's action IDs, Julia's source hashes, Lean's Lake traces and
Taichi's kernel hashes are all hashes of named or textual content
{[}K{]}. The Wasm binary comes closest, because names sit in an optional
custom section {[}K{]}. So R has to become an operation the framework
applies to any IR (rename bound names to positions, strip paths, then
hash), not a property compilers are claimed to have. The project's
experiment E5 implements that operation in two forms and measures them
on a small corpus (§8.3).

\textbf{L fails as a description; the bug counts are the evidence.}
Csmith found more than 325 previously unknown bugs in C compilers
{[}1{]}. EMI found 147 confirmed bugs in its first 11 months and 1,602
GCC and LLVM bugs in total, 1,007 of them fixed {[}2{]}. Alive2, a
bounded translation validator for LLVM, found 47 new bugs (28 fixed) and
prompted 8 patches to the LLVM language reference; a follow-up encoding
of LLVM's memory model found 21 more {[}15{]}. Zig's self-hosted x86-64
backend, the default for Debug builds since 0.15, passes 1,984 of 2,008
behaviour tests (the project's development log gives 1,987 against 1,980
for LLVM; the two sources disagree) {[}3{]}, so lowering is not total in
a shipping compiler. CompCert and CakeML show that L can hold as a
theorem: CakeML carries its guarantee down to machine code for five
instruction sets and has compiled itself inside the HOL4 logic {[}16{]}.
Neither is a production compiler for the languages surveyed. L survives
only as a norm: a step either produces output its checker accepts or
stops, and production compilers are measured by how far they fall short.

\textbf{C is present but local.} Every compiler has cost models, but
almost all are local heuristics: inline budgets, LLVM's target cost
queries and scheduling models, GCC's cost hooks {[}K{]}. Cranelift's
acyclic e-graph is the only production case found that ranks the members
of an explicit equivalence class and extracts the cheapest, and it
deliberately saturates once, within a budget {[}17{]}. Research equality
saturation shows uneven gains: TENSAT reports up to 16\% {[}18{]}, and
equality saturation inside XLA averages 3.45\% with a 56.26\% peak on
one model {[}19{]}. Only one mechanised result about a pass's
\emph{performance} guarantee was found, for an inlining heuristic in
Rocq {[}20{]}. Every cost found is in cycles, latency, bytes or node
counts; none is in joules, and the shortest program is not the cheapest
whenever unrolling, vectorising or table lookup makes code longer but
faster {[}K{]}.

\textbf{K is the strongest part, but it guards the source more than the
lowering.} Type systems, Rust's borrow checker and new trait solver
{[}21{]}, GHC's Core Lint, LFortran's ASR verifier, MLIR op verifiers,
Futhark's typed IR, Lean's kernel, SPARK's GNATprove proof levels
{[}22{]} and Erlang's code validator {[}K{]} are all real checkers.
Three qualifications matter. Checking of the lowering itself exists only
in pieces: Alive2 for some LLVM peepholes, SMT verification of some
Cranelift ISLE rules, the mechanised Wasm specification {[}K{]}. Many IR
checkers are debug-only: Core Lint is opt-in and the LLVM verifier
between passes is commonly off in release builds {[}K{]}. And bounded
checkers can miss bugs: Alive2 unrolls loops to a bound {[}15{]}, and
EMI and Csmith find bugs but cannot certify their absence. The ML
compilers' checks are tolerance tests on sample inputs, which are tests,
not certificates.

\textbf{U is the part most often present, but its keys hash names and
paths.} Go's build cache maps an action ID (a hash of every input to a
step) to an output {[}23{]}; rustc uses incremental fingerprints; LLVM's
ThinLTO cache, Wasmtime's module cache, JAX's compilation cache,
Triton's kernel cache and TVM's tuning database are all keyed by hashes
{[}K{]}. Reuse of verified work compounds trust: CakeML compiled its own
verified binary inside the logic {[}16{]}; ``Compiling by Proving''
turns proofs into optimised rewrite rules {[}24{]}; and DDC counters the
trusting-trust attack under stated assumptions {[}25{]}. Two
qualifications remain. The keys hash names and paths, so equal programs
with different names miss the cache. And TensorRT-LLM abandoned its
serialised, verified-once engines for PyTorch as its only backend
{[}26{]}: stored artefacts lost to a flexible eager path when hardware
and models changed faster than the artefacts could be rebuilt.

\textbf{E is absent everywhere.} No compiler in any of the twelve
families accounts for energy {[}K; NOT FOUND for any counterexample{]}.
Energy is measured outside compilers: by machine (the June 2026 Green500
leader runs at 73.28 GFlops per watt {[}27{]}) and by counters such as
RAPL and NVML {[}K{]}. Research on energy-aware compilation is modest:
the choice of compiler flags shifts energy by tens of percent, and
energy and run time are ``roughly correlated, though not always''
{[}28{]}; a 2025 Fortran study found ``minimal difference'' in energy
between gfortran, ifx and ifort, tracking run time {[}29{]}. The only
compiler found that treats energy as a first-class quantity is the
research language Eel, which treats reversible operations as free and
charges for bit creation and destruction {[}30{]}. Either E is not a
necessary part of a law of compilation, or the whole field falls short
of it. This study takes the second reading: E is a part the framework
adds, supplied by an external meter.

\textbf{G is present only in some compilers.} rustc emits no code after
a type or borrow error; Go refuses unused imports and variables; Wasm
engines refuse an invalid module at load time; Erlang's validator checks
generated code before writing it {[}K{]}. SPARK is a gate only when a
project's CI requires GNATprove with no unproved checks; the compiler
itself does not refuse {[}K{]}. Elsewhere G is off or partial. tsc emits
JavaScript despite type errors unless \texttt{noEmitOnError} is set;
LLVM and GCC ship from trunk on post-commit testing and reverts rather
than a gate before merging; Zig analyses lazily, so code that is never
referenced is never checked {[}K{]}. ML compilers have fallbacks rather
than gates: a graph break in TorchDynamo falls back to eager Python
unless \texttt{fullgraph=True} turns it into an error, and ONNX Runtime
sends unsupported nodes to its CPU provider {[}K{]}.

\section{4. The convergence function}\label{the-convergence-function}

\subsection{4.1 Landauer and Kolmogorov as conceptual
bounds}\label{landauer-and-kolmogorov-as-conceptual-bounds}

The candidate's convergence function says that the cost of a compiled
program has a floor at Landauer's limit and a target at
Kolmogorov-minimal size. \textbf{Landauer's limit} is the minimum heat,
\(k_B T \ln 2\), released into surroundings at temperature \(T\) when
one bit is erased: about 2.871 zeptojoules at 300 K (2.853 at 298.15 K)
{[}K, arithmetic{]}. It bounds heat per bit \emph{erased}, not per
operation, and it is reached only in the infinitely slow limit; erasing
faster costs more, roughly as \(1/\tau\) for erasure time \(\tau\)
{[}K{]}. Single-bit experiments confirm the bound at that limit, in
colloidal beads (Bérut et al., 2012) and in nanomagnetic memory (Hong et
al., 2016) {[}31{]}. \textbf{Kolmogorov complexity} \(K(x)\) is the
length of the shortest program that prints \(x\). Both are real physics
or mathematics, and as bounds they say something true: no compiled
program can dissipate less than its erasures require, and none can be
shorter than the shortest equivalent program.

\subsection{4.2 Why neither is computable or measurable as
stated}\label{why-neither-is-computable-or-measurable-as-stated}

\textbf{The Kolmogorov target cannot be computed.} \(K(x)\) is
uncomputable (Kolmogorov, Chaitin, Solomonoff); compressors give only
upper bounds, and program equivalence, the constraint the target sits
under, is undecidable in general by Rice's theorem {[}K{]}. The
full-employment theorem for compiler writers makes this concrete: no
compiler can always produce the smallest equivalent program, because
that would decide halting {[}K{]}. Superoptimisers approach the target
for short loop-free fragments under an SMT or test-vector check, but
never with a guarantee.

\textbf{The Landauer floor cannot be computed either, and it is far
below today's costs.} Computing the optimal heat functions for one fixed
circuit is PP-complete according to the abstract of a 2026 paper,
multiplicative approximation is NP-hard, and only additive approximation
has a polynomial-time randomised algorithm {[}4{]}; the body reportedly
argues PP-hardness, so membership in PP remains to be checked. Choosing
among equivalent circuits, which is the compilation problem, sits on top
of that. At the commonly quoted 45 nm figures, a 32-bit integer add
costs about 0.1 picojoules, \textbf{about \(3.5\times10^{7}\) times}
\(k_B T \ln 2\), or about \(10^6\) per erased bit if 32 bits are
overwritten; a 640-picojoule DRAM read is \textbf{about
\(2.2\times10^{11}\)} times the floor; and the overhead of one
out-of-order instruction, about 70 picojoules, is \textbf{about
\(2.5\times10^{10}\)} {[}32{]}. Today's gap is therefore \(10^6\) to
\(10^{11}\), and almost all of it is wires, memory and control, not
logical erasure.

\textbf{The two terms overlap.} Bennett showed that any computation can
be made logically reversible, so the Landauer floor applies only to
information a program deliberately discards {[}K{]}. Zurek, and Bennett
and colleagues, showed that the minimum thermodynamic cost of erasing a
string \(x\) is about \(k_B T \ln 2 \cdot K(x)\), and that converting
\(x\) to \(y\) costs about their information distance {[}K{]}. The
Landauer cost of erasure is already a Kolmogorov quantity, so minimising
both jointly counts the same thing twice. Three further problems remain.
Shortness and low energy trade off, so a joint objective without a
stated weighting is a Pareto front, not a function. Time is missing from
the statement. And the energy spent compiling can outweigh the savings:
in an October 2026 study, 42 of 44 configurations of model
self-refinement used 4\% to 886\% more execution energy than the
reference solutions, and the two exceptions broke even only after about
\(10^6\) to \(10^7\) executions {[}33{]}.

One result supports the idea at its core. Wolpert and Kolchinsky's
``circuit Landauer loss'' shows that two circuits computing the same
function can have different minimum dissipation {[}34{]}, and
mismatch-cost results give strictly positive lower bounds for circuits
{[}35{]}. Circuit structure, which is what a compiler chooses, does
change the physical floor.

\subsection{4.3 The measurable
replacement}\label{the-measurable-replacement}

The replacement keeps both terms and makes each computable:

\begin{enumerate}
\def\labelenumi{\arabic{enumi}.}
\tightlist
\item
  \textbf{ε, the erasure count}: the number of bits a run overwrites
  irreversibly, computed from the IR (bits overwritten or discarded
  without having been copied out). \(\varepsilon \cdot k_B T \ln 2\) is
  the run's Landauer bound.
\item
  \textbf{λ, a declared description-length proxy}: IR node count, binary
  bytes, compressed bytes, or an MDL code length relative to a named
  model class. The proxy is declared before a measurement and not
  changed after it.
\end{enumerate}

Both are computed under an equivalence checked on a decidable fragment,
and reported next to \textbf{measured joules} in the energy regime the
hardware belongs to:

{\footnotesize\setlength{\tabcolsep}{3.8pt}\renewcommand{\arraystretch}{1.15}
\begin{longtable}[]{@{}
  >{\raggedright\arraybackslash}p{(\linewidth - 4\tabcolsep) * \real{0.1146}}
  >{\raggedright\arraybackslash}p{(\linewidth - 4\tabcolsep) * \real{0.2708}}
  >{\raggedright\arraybackslash}p{(\linewidth - 4\tabcolsep) * \real{0.6146}}@{}}
\toprule\noalign{}
\begin{minipage}[b]{\linewidth}\raggedright
\textbf{Regime}
\end{minipage} & \begin{minipage}[b]{\linewidth}\raggedright
\textbf{Energy is proportional to}
\end{minipage} & \begin{minipage}[b]{\linewidth}\raggedright
\textbf{Example hardware}
\end{minipage} \\
\midrule\noalign{}
\endhead
\bottomrule\noalign{}
\endlastfoot
A & Counted operations & In-order cores, FPGAs \\
B & Time multiplied by power & Out-of-order CPUs, GPUs \\
C & Entropy production & Thermodynamic hardware, or a simulation of it;
the only regime in which Landauer units apply directly \\
\end{longtable}
}

Regimes A and B are approximations: static power scales with time even
on in-order cores {[}K{]}. The claim becomes testable: does lowering ε
and λ predict lower measured energy in each regime? No published
measurement of the erasure content of real compiled programs against
their energy was found {[}NOT FOUND{]}, so the result would be new
evidence whichever way it comes out. Self-published preprints that claim
convergence to the Landauer limit at a rate of \(O(1/\sqrt{N})\) are
unreviewed, and their rate is a statistic about estimating work, not a
rate at which programs approach a floor {[}36{]}; they are not cited as
support.

\section{5. The gear ladder and the 98/2
split}\label{the-gear-ladder-and-the-982-split}

RCE Compile orders its own compilation steps as a \textbf{gear ladder}
of six rungs, G0 to G5, tried in order of cost. G0 reuses a stored
result by its hash; G1 applies rewrite rules proved correct in advance;
G2 extracts the cheapest form from an e-graph. These three are
deterministic. G3 to G5 are stochastic proposers, and every output they
produce must pass the gate before it is released; the charter defines G3
to G5 only as a group, and their separate definitions are open {[}6{]}.
(The rungs are written G0 to G5; the Form's gate is the part named G.)
The design principle attached to the ladder is a split: \textbf{at least
98\% of the work deterministic, at most 2\% stochastic, and everything
gated.}

No source measures such a split for any compiler, by decision count, CPU
time, cost or energy {[}NOT FOUND{]}. The verified systems sit at the
two ends instead. Equality saturation, Souper, Cranelift and MLGO are
fully deterministic at compile time; MLGO's policy is fixed and compiled
into LLVM, and it decides one class of choices (inline or not, which
register to evict) for gains of 0.3\% to 7\% {[}37{]}. At the other end,
a stochastic author writes nearly everything behind a deterministic
gate: in work reported in October 2026, a model wrote about 180 million
lines of Rocq proof that the Rocq checker accepts or rejects in full
{[}38{]}, and AdaCore's GNAT Foundry: Intersection demonstrator has an
agent write SPARK while deterministic tools produce all the release
evidence {[}39{]}. By authoring effort that split is inverted, while the
gate is total, so any claim about 98/2 must first say what is being
counted. One taxonomy places hybrid workflows at 14\% to 56\%
deterministic, but it is a classification, not a measurement {[}40{]}.

The literature does support the two principles beneath the split.
\textbf{Gate everything}: SuperCoder's best configuration still fails
5\% of its test-checked outputs, and the strongest commercial baseline
in that paper passes only 51.5\% {[}41{]}; STOKE paired a stochastic
search with an SMT validator {[}K{]}. \textbf{Put deterministic reuse
first}: Compiled AI moves all stochastic work into a one-time compile
phase and breaks even after about 17 runs {[}42{]}, and Progressive
Crystallization promotes recurring agent runs into deterministic
workflows {[}43{]}. One caution carries over: in Veritas and
CktFormalizer the gate checks the stochastic output against its own
specification (CNF clauses, hardware types), while an error in the
specification sits outside the gate {[}44{]} {[}45{]}.

The 98/2 split is therefore a registered hypothesis, listed in §7, not a
measured result. It is to be measured rung by rung with a declared unit
and recorded by the gate. Candidate units are accepted decisions,
accepted rewrites, CPU seconds and joules; each is reported separately,
because they can disagree.

\section{6. Prior art}\label{prior-art}

No earlier source found names a single law of compilation (§2.1). The
closest earlier claims each cover a few of the seven parts.

\begin{itemize}
\tightlist
\item
  \textbf{Morris (1973)}, following Burstall and Landin (1969) and
  refined by the ADJ group (Thatcher, Wagner and Wright, 1979 to 1981),
  made compiler correctness a commuting square of algebra homomorphisms:
  compiling then interpreting equals interpreting then encoding {[}46{]}
  {[}47{]}. That is L and K. Lowering as stated in the Form restates it
  and claims no novelty there.
\item
  \textbf{Futamura (1971)} showed that specialising an interpreter on a
  program yields a compiled program, that specialising the specialiser
  on the interpreter yields a compiler, and that specialising it on
  itself yields a compiler generator {[}48{]}. \textbf{Ershov's}
  Novosibirsk school reached the same view as ``mixed computation''
  {[}49{]}. That is L and U, and the strongest prior claim that
  compilers are derived rather than designed.
\item
  \textbf{Hoare, He and Sampaio (Acta Informatica, 1993)} treated
  compilation as reducing a source program, by algebraic laws of
  programming, to a normal form that mirrors the target machine, with
  correctness following from the laws; the work fed into Hoare and He's
  \emph{Unifying Theories of Programming} (1998) {[}50{]}. This is the
  closest prior work to calling compilation ``governed by laws'', and
  the closest to the normal-form framing this study adopts.
\item
  \textbf{Cousot and Cousot (POPL 1977)} unified static analyses as
  fixpoints over lattices, and later extended this to program
  transformations {[}51{]}. That is K, and the precision ordering in C.
\item
  \textbf{Eel (Tyagi, Lynch and Demaine, MIT, 2016)} is the only
  compiler found that ties code to an energy floor, by charging for bit
  creation and destruction {[}30{]}. That is E.
\end{itemize}

UNCOL (1958) tried and failed to give one universal R; MLIR answers that
failure with many levels but holds no semantics {[}52{]}. E-graphs
already combine canonical sets, cost-ranked extraction and sound
rewrites in one structure {[}53{]}.

\textbf{Two primary sources have not yet been read in full: Hoare, He
and Sampaio (1993) {[}50{]} and Ershov's ``On the essence of
compilation'' (1977) {[}54{]}.} Both are cited from memory {[}K{]}.
Either may contain a precedent for the framing or the wording, and both
must be read before RCE Compile claims priority for anything. The novel
part that remains, on present evidence, is the combination: one
statement binding a physical floor and an information-theoretic target
to verified, reuse-compounding lowering, with energy as a ranked
quantity.

\section{7. Hypotheses and falsifiers}\label{hypotheses-and-falsifiers}

The project registered seven hypotheses on 1 November 2026, before any
measurement of its own, each with a metric, a condition for support, a
condition for falsification and an experiment {[}6{]}. The statements
and falsifiers below are quoted from the registration. A \emph{declared
corpus} is a set of programs fixed, by path and hash, before the
measurement that uses it. A change after registration is recorded as a
dated row in the registration's table of changes. Besides the
registration itself, that table holds five changes, all dated 2 November
2026: a verdict on the one E1 difference in Go that survived remapping
on Linux, which adds the import path to H1's declared inputs (§8.1); a
record that the first macOS run of E1 has no verdict yet (§8.1); a
change to the metrics of H2 and H3, which now record the authority of
each gate's evidence and the digests its verdict is bound to (§3.3,
§8.2), made before either experiment has run; the attachment of
experiment E5, with its predictions and falsifiers, to the fifth
hypothesis below, before E5's first run (§8.3); and E5's first results,
which support that hypothesis on E4's corpus by its registered
conditions (§8.3).

{\footnotesize\setlength{\tabcolsep}{3.8pt}\renewcommand{\arraystretch}{1.15}
\begin{longtable}[]{@{}
  >{\raggedright\arraybackslash}p{(\linewidth - 6\tabcolsep) * \real{0.0387}}
  >{\raggedright\arraybackslash}p{(\linewidth - 6\tabcolsep) * \real{0.3355}}
  >{\raggedright\arraybackslash}p{(\linewidth - 6\tabcolsep) * \real{0.3355}}
  >{\raggedright\arraybackslash}p{(\linewidth - 6\tabcolsep) * \real{0.2903}}@{}}
\toprule\noalign{}
\begin{minipage}[b]{\linewidth}\raggedright
\end{minipage} & \begin{minipage}[b]{\linewidth}\raggedright
\textbf{Statement}
\end{minipage} & \begin{minipage}[b]{\linewidth}\raggedright
\textbf{Falsified if}
\end{minipage} & \begin{minipage}[b]{\linewidth}\raggedright
\textbf{Experiment}
\end{minipage} \\
\midrule\noalign{}
\endhead
\bottomrule\noalign{}
\endlastfoot
H1 & Lowering is a function of declared inputs. With declared inputs
(source, flags, toolchain, remapped build path, fixed
\texttt{SOURCE\_\allowbreak{}DATE\_\allowbreak{}EPOCH}), each compiler's output bytes are
identical across directories and across machines. & Any byte difference
survives remapping, across directories or across machines. & E1, run
(§8.1) \\
H2 & Gates can be classified. Gates differ by compiler and can be
classified: some emit nothing by default when a check fails, others emit
unless configured not to. & A compiler's documented gate behaviour
differs from what the injected-fault test observes. & E2, protocol
(§8.2) \\
H3 & Every pass and flag has an equivalence level. Every optimisation
pass and flag can be assigned one of the four equivalence levels. & A
pass's observed output change exceeds its assigned level. & E3, protocol
(§8.2) \\
H4 & Erasure and description length against measured energy. The counted
irreversible erasure ε plus a declared description-length proxy λ
predict measured energy in regime C (or a simulation of it), but explain
less variance than operation counts in regime A and than time multiplied
by power in regime B. & ε and λ explain no variance in regime C, or
explain more variance than the regime's own baseline in A or B, which
would contradict the scale-gap analysis of the charter. & E4, erasure
counts + Regime-B Linux RAPL \texttt{reported\_\allowbreak{}j} \& payback horizons on
20 kernels (§8.2); Regimes A and C open \\
H5 & A name-free hash raises reuse without false hits. A name-free hash
(bound names replaced by positions, path removed) raises reuse hits on
renamed but equal programs, without false hits. & There is any false
hit, or no gain in hits on the declared corpus. & E5, supported on E4's
corpus and scaled to 245 programs + 500 verified semantic mutants
(§8.3) \\
H6 & Deterministic rungs decide at least 98\%. The deterministic rungs
G0 to G2 decide at least 98\% of accepted decisions on a declared
corpus, and every stochastic proposal is gated. & The measured
stochastic share is above 2\% in the declared unit, or any stochastic
output is released without a recorded gate decision. & E6, planned;
needs a compiler built on the gear ladder, which does not exist yet \\
H7 & Partial lowering with stops admits no miscompilation inside its
fragment. Lowering stated as partial, with a stop as an outcome, has a
measurable stop rate, and gated lowering admits no miscompilation that
EMI-style testing finds inside the checker's decidable fragment. &
EMI-style testing finds a miscompilation inside the fragment the gate
claims to cover. & E7, planned; no protocol yet \\
\end{longtable}
}

The conditions for support are the converse of each falsifier, with
these specifics from the registration: H1 is supported if the
differing-byte count is 0 for every output file in every pair; H2 if
every compiler falls into one class and its observed behaviour matches
its documented behaviour; H4 if ε and λ explain variance in regime C
above a threshold declared in the E4 protocol, and less than the
baseline in regimes A and B; H6 if the stochastic share is at most 2\%
in the declared unit and every stochastic output has a recorded gate
decision; H7 if the stop rate is measured and no such miscompilation is
found.

\section{8. Experiments}\label{experiments}

\subsection{8.1 E1: determinism of lowering,
measured}\label{e1-determinism-of-lowering-measured}

\textbf{Method.} E1 tests H1 {[}6{]}. Each \emph{configuration} is one
compiler with one flag set, and names one or two small programs. For
every program, the harness makes four pairs of builds, each build in a
fresh directory that holds only the copied source: the same directory
twice; two directories with absolute paths of different lengths; the
same two directories with the toolchain's documented path remapping; and
the same directory twice with \texttt{SOURCE\_\allowbreak{}DATE\_\allowbreak{}EPOCH} unset, at
least 1.1 s apart. All builds but the last pair run with a declared
environment (\texttt{LC\_ALL=C}, \texttt{TZ=UTC},
\texttt{SOURCE\_\allowbreak{}DATE\_\allowbreak{}EPOCH=1735689600}), and every build gets new,
empty compiler caches, so the second build of a pair is a real rebuild.
For each pair the harness records, per output file, the number of byte
positions at which the two files differ, and, for an ELF file, how many
of those bytes fall in each section. The count is positional, so a path
of another length that shifts the rest of a file counts every shifted
byte. The remappings are \texttt{-\allowbreak{}ffile-\allowbreak{}prefix-\allowbreak{}map} for gcc, g++, clang,
clang++, gfortran and gnatmake, \texttt{-\allowbreak{}\/\allowbreak{}-\allowbreak{}remap-\allowbreak{}path-\allowbreak{}prefix} for
rustc, \texttt{-trimpath} for Go, \texttt{+deterministic} for erlc and
\texttt{BUILD\_\allowbreak{}PATH\_\allowbreak{}PREFIX\_\allowbreak{}MAP} for ocamlopt; Go 1.21 was described by
its developers as ``the first Go toolchain with perfectly reproducible
builds'' {[}55{]}. Julia has no remapping in the matrix: its remapped
run writes the code without source lines. No remapping was found for
Lean, \texttt{wasmtime\ compile}, Zig, GHC or GnuCOBOL, so their
remapped comparison is not tested. \texttt{go\ build} is measured three
times: in module mode, Go's default; in GOPATH mode with the source file
named on the command line; and in GOPATH mode with the source placed in
a \texttt{GOPATH} and built by its import path, \texttt{e1}. For each Go
configuration the harness runs the remapped build once more in each
directory with \texttt{GODEBUG=gocachehash=1}, which makes
\texttt{go\ build} print every input it hashes into each action ID, and
records the lines that name the build directory. The harness also
records the SHA-256 of every output, so that runs on two machines can be
compared where a tool's version string is the same on both. Results are
written by the harness, and the write-up is generated from them by a
script; no figure in it is typed.

\textbf{Machines and toolchains.} Three runs on 2 November 2026 (UTC).
Two are on Linux x86-64 with 4 CPUs and Ubuntu 24.04.5. The \emph{local
machine} (Intel Xeon at 2.80 GHz) has gcc, g++ and gfortran 13.3.0,
clang and clang++ 18.1.3, rustc 1.97.0 (including its
\texttt{wasm32-\allowbreak{}wasip1} target), go 1.24.7, Lean 4.12.0 and Wasmtime
26.0.1. The \emph{CI runner}, a GitHub-hosted runner (Intel Xeon
Platinum 8573C) on which a workflow in the repository runs the same
harness, has gcc, g++, gfortran and gnatmake 13.3.0, clang and clang++
18.1.3, rustc 1.99.1, go 1.24.13, Zig 0.17.0, GHC 9.14.1, GnuCOBOL
3.1.2.0, OCaml 4.14.1, Erlang/OTP 25 and Julia 1.13.2; Lean and Wasmtime
are not installed there. The third, the \emph{macOS laptop}, is recorded
in the repository's table of changes as a MacBook Pro (M5 Max) with
macOS 27.2 on arm64; its results file records 18 CPUs and neither the
CPU model nor the system release. On it \texttt{gcc}, \texttt{g++},
\texttt{clang} and \texttt{clang++} all run Apple clang 21.0.0, beside
gfortran 16.2.0 (Homebrew), rustc 1.99.1 (Homebrew), go 1.27.3, Lean
4.35.0, Wasmtime 50.0.1, Zig 0.17.0, GHC 9.14.1, GnuCOBOL 3.2.0, OCaml
5.5.2 and Julia 1.13.3; \texttt{gnatmake}, \texttt{erlc} and rustc's
WebAssembly target are not installed there. Its executables and object
files are Mach-O files, not ELF, and a \texttt{-g} build of C, C++,
Fortran or Rust leaves its debug information in a separate
\texttt{.dSYM} bundle beside the executable. The file was written before
the harness could attribute differing Mach-O bytes to regions of the
file, so for this run E1 records how many bytes differ but not where.

\textbf{Results, local machine.} Fifteen configurations on 28 programs
in all (two per configuration, one for Lean and one for Wasmtime).

{\footnotesize\setlength{\tabcolsep}{3.8pt}\renewcommand{\arraystretch}{1.15}
\begin{longtable}[]{@{}
  >{\raggedright\arraybackslash}p{(\linewidth - 8\tabcolsep) * \real{0.2095}}
  >{\raggedright\arraybackslash}p{(\linewidth - 8\tabcolsep) * \real{0.1714}}
  >{\raggedright\arraybackslash}p{(\linewidth - 8\tabcolsep) * \real{0.2190}}
  >{\raggedright\arraybackslash}p{(\linewidth - 8\tabcolsep) * \real{0.2000}}
  >{\raggedright\arraybackslash}p{(\linewidth - 8\tabcolsep) * \real{0.2000}}@{}}
\toprule\noalign{}
\begin{minipage}[b]{\linewidth}\raggedright
\textbf{Configuration}
\end{minipage} & \begin{minipage}[b]{\linewidth}\raggedright
\textbf{Same directory twice}
\end{minipage} & \begin{minipage}[b]{\linewidth}\raggedright
\textbf{Two directories}
\end{minipage} & \begin{minipage}[b]{\linewidth}\raggedright
\textbf{Two directories, remapped}
\end{minipage} & \begin{minipage}[b]{\linewidth}\raggedright
\textbf{Same directory, clock free}
\end{minipage} \\
\midrule\noalign{}
\endhead
\bottomrule\noalign{}
\endlastfoot
\texttt{gcc\ -\allowbreak{}O2\ -\allowbreak{}g} (C) & identical (2 of 2) & differs (up to 2,131
bytes) & identical & identical \\
\texttt{gcc\ -O2} (C) & identical & identical & identical & identical \\
\texttt{clang\ -\allowbreak{}O2\ -\allowbreak{}g} (C) & identical & differs (up to 2,191 bytes) &
identical & identical \\
\texttt{g++\ -\allowbreak{}O2\ -\allowbreak{}g} (C++) & identical & differs (up to 20,602 bytes) &
identical & identical \\
\texttt{clang++\ -\allowbreak{}O2\ -\allowbreak{}g} (C++) & identical & differs (up to 91,041
bytes) & identical & identical \\
\texttt{rustc\ -\allowbreak{}O\ -\allowbreak{}g} (Rust) & identical & differs (up to 2,441,330
bytes) & identical & identical \\
\texttt{rustc\ -O} (Rust) & identical & identical & identical &
identical \\
\texttt{rustc} to WebAssembly, \texttt{-O\ -g} & identical & differs (up
to 1,213,811 bytes) & identical & identical \\
\texttt{go\ build} (module mode) & identical & differs (up to 741,728
bytes) & identical & identical \\
\texttt{go\ build}, GOPATH mode & identical & differs (up to 743,688
bytes) & \textbf{differs (58 and 59 bytes)} & identical \\
\texttt{go\ build}, GOPATH mode, import path declared & identical &
differs (up to 1,058,112 bytes) & identical & identical \\
\texttt{gfortran\ -\allowbreak{}O2\ -\allowbreak{}g} (Fortran) & identical & differs (up to 2,368
bytes) & identical & identical \\
\texttt{gfortran\ -\allowbreak{}O2} (Fortran) & identical & identical & identical &
identical \\
\texttt{lean\ -c} (Lean 4, 1 program) & identical & identical & not
tested & identical \\
\texttt{wasmtime\ compile} (Cranelift, 1 program) & identical &
identical & not tested & identical \\
\end{longtable}
}

\textbf{Results, CI runner.} Twenty-one configurations on 34 programs in
all.

{\footnotesize\setlength{\tabcolsep}{3.8pt}\renewcommand{\arraystretch}{1.15}
\begin{longtable}[]{@{}
  >{\raggedright\arraybackslash}p{(\linewidth - 8\tabcolsep) * \real{0.2170}}
  >{\raggedright\arraybackslash}p{(\linewidth - 8\tabcolsep) * \real{0.1698}}
  >{\raggedright\arraybackslash}p{(\linewidth - 8\tabcolsep) * \real{0.2170}}
  >{\raggedright\arraybackslash}p{(\linewidth - 8\tabcolsep) * \real{0.1981}}
  >{\raggedright\arraybackslash}p{(\linewidth - 8\tabcolsep) * \real{0.1981}}@{}}
\toprule\noalign{}
\begin{minipage}[b]{\linewidth}\raggedright
\textbf{Configuration}
\end{minipage} & \begin{minipage}[b]{\linewidth}\raggedright
\textbf{Same directory twice}
\end{minipage} & \begin{minipage}[b]{\linewidth}\raggedright
\textbf{Two directories}
\end{minipage} & \begin{minipage}[b]{\linewidth}\raggedright
\textbf{Two directories, remapped}
\end{minipage} & \begin{minipage}[b]{\linewidth}\raggedright
\textbf{Same directory, clock free}
\end{minipage} \\
\midrule\noalign{}
\endhead
\bottomrule\noalign{}
\endlastfoot
\texttt{gcc\ -\allowbreak{}O2\ -\allowbreak{}g} (C) & identical (2 of 2) & differs (up to 2,159
bytes) & identical & identical \\
\texttt{gcc\ -O2} (C) & identical & identical & identical & identical \\
\texttt{clang\ -\allowbreak{}O2\ -\allowbreak{}g} (C) & identical & differs (up to 2,129 bytes) &
identical & identical \\
\texttt{g++\ -\allowbreak{}O2\ -\allowbreak{}g} (C++) & identical & differs (up to 20,635 bytes) &
identical & identical \\
\texttt{clang++\ -\allowbreak{}O2\ -\allowbreak{}g} (C++) & identical & differs (up to 91,003
bytes) & identical & identical \\
\texttt{rustc\ -\allowbreak{}O\ -\allowbreak{}g} (Rust) & identical & differs (up to 2,449,365
bytes) & identical & identical \\
\texttt{rustc\ -O} (Rust) & identical & identical & identical &
identical \\
\texttt{rustc} to WebAssembly, \texttt{-O\ -g} & identical & differs (up
to 1,249,142 bytes) & identical & identical \\
\texttt{go\ build} (module mode) & identical & differs (up to 1,061,816
bytes) & identical & identical \\
\texttt{go\ build}, GOPATH mode & identical & differs (up to 1,061,871
bytes) & \textbf{differs (59 and 60 bytes)} & identical \\
\texttt{go\ build}, GOPATH mode, import path declared & identical &
differs (up to 751,516 bytes) & identical & identical \\
\texttt{gfortran\ -\allowbreak{}O2\ -\allowbreak{}g} (Fortran) & identical & differs (up to 2,403
bytes) & identical & identical \\
\texttt{gfortran\ -\allowbreak{}O2} (Fortran) & identical & identical & identical &
identical \\
\texttt{zig\ build-\allowbreak{}exe} (Zig, 1 program) & \textbf{differs (1,488,011
bytes)} & differs (3,790,312 bytes) & not tested & differs (1,858,521
bytes) \\
\texttt{zig\ build-\allowbreak{}exe\ -\allowbreak{}O\ ReleaseSmall} (Zig, 1 program) & identical &
identical & not tested & identical \\
\texttt{ghc\ -O1} (Haskell, 1 program) & identical & identical & not
tested & identical \\
\texttt{gnatmake\ -\allowbreak{}O2\ -\allowbreak{}g} (Ada, 1 program) & identical & differs
(26,349 bytes) & \textbf{differs (12,129 bytes)} & differs (3 bytes) \\
\texttt{cobc\ -x} (COBOL, 1 program) & \textbf{differs (28 bytes)} &
differs (11,891 bytes) & not tested & differs (25,886 bytes) \\
\texttt{ocamlopt\ -\allowbreak{}g} (OCaml, 1 program) & identical & differs (164,436
bytes) & identical & identical \\
\texttt{erlc} (Erlang, 1 program) & identical & differs (197 bytes) &
identical & identical \\
Julia \texttt{code\_\allowbreak{}llvm} and \texttt{code\_\allowbreak{}native} (1 program) &
\textbf{differs (168 bytes)} & differs (144,945 bytes) & \textbf{differs
(146 bytes)} & differs (168 bytes) \\
\end{longtable}
}

\textbf{Results, macOS laptop.} Twenty configurations on 32 programs in
all; \texttt{gcc} and \texttt{g++} are Apple clang here.

{\footnotesize\setlength{\tabcolsep}{3.8pt}\renewcommand{\arraystretch}{1.15}
\begin{longtable}[]{@{}
  >{\raggedright\arraybackslash}p{(\linewidth - 8\tabcolsep) * \real{0.2110}}
  >{\raggedright\arraybackslash}p{(\linewidth - 8\tabcolsep) * \real{0.1835}}
  >{\raggedright\arraybackslash}p{(\linewidth - 8\tabcolsep) * \real{0.2110}}
  >{\raggedright\arraybackslash}p{(\linewidth - 8\tabcolsep) * \real{0.2018}}
  >{\raggedright\arraybackslash}p{(\linewidth - 8\tabcolsep) * \real{0.1927}}@{}}
\toprule\noalign{}
\begin{minipage}[b]{\linewidth}\raggedright
\textbf{Configuration}
\end{minipage} & \begin{minipage}[b]{\linewidth}\raggedright
\textbf{Same directory twice}
\end{minipage} & \begin{minipage}[b]{\linewidth}\raggedright
\textbf{Two directories}
\end{minipage} & \begin{minipage}[b]{\linewidth}\raggedright
\textbf{Two directories, remapped}
\end{minipage} & \begin{minipage}[b]{\linewidth}\raggedright
\textbf{Same directory, clock free}
\end{minipage} \\
\midrule\noalign{}
\endhead
\bottomrule\noalign{}
\endlastfoot
\texttt{gcc\ -\allowbreak{}O2\ -\allowbreak{}g} (C) & \textbf{differs (up to 101 bytes)} & differs
(up to 1,317 bytes) & \textbf{differs (up to 1,316 bytes)} & differs (up
to 103 bytes) \\
\texttt{gcc\ -O2} (C) & identical (2 of 2) & identical & identical &
identical \\
\texttt{clang\ -\allowbreak{}O2\ -\allowbreak{}g} (C) & \textbf{differs (up to 38 bytes)} &
differs (up to 1,318 bytes) & \textbf{differs (up to 1,319 bytes)} &
differs (up to 101 bytes) \\
\texttt{g++\ -\allowbreak{}O2\ -\allowbreak{}g} (C++) & \textbf{differs (up to 38 bytes)} &
differs (up to 126,164 bytes) & \textbf{differs (up to 126,165 bytes)} &
differs (up to 134 bytes) \\
\texttt{clang++\ -\allowbreak{}O2\ -\allowbreak{}g} (C++) & \textbf{differs (up to 38 bytes)} &
differs (up to 126,165 bytes) & \textbf{differs (up to 126,164 bytes)} &
differs (up to 135 bytes) \\
\texttt{rustc\ -\allowbreak{}O\ -\allowbreak{}g} (Rust) & identical & differs (up to 105,506
bytes) & \textbf{differs (up to 105,506 bytes)} & identical \\
\texttt{rustc\ -O} (Rust) & identical & identical & identical &
identical \\
\texttt{go\ build} (module mode) & identical & differs (up to 1,089,596
bytes) & identical & identical \\
\texttt{go\ build}, GOPATH mode & identical & differs (up to 1,089,588
bytes) & \textbf{differs (119 and 120 bytes)} & identical \\
\texttt{go\ build}, GOPATH mode, import path declared & identical &
differs (up to 1,089,997 bytes) & identical & identical \\
\texttt{gfortran\ -\allowbreak{}O2\ -\allowbreak{}g} (Fortran) & \textbf{differs (up to 38 bytes)}
& differs (up to 1,015 bytes) & \textbf{differs (up to 1,014 bytes)} &
differs (up to 103 bytes) \\
\texttt{gfortran\ -\allowbreak{}O2} (Fortran) & identical & identical & identical &
identical \\
\texttt{lean\ -c} (Lean 4, 1 program) & identical & identical & not
tested & identical \\
\texttt{wasmtime\ compile} (Cranelift, 1 program) & identical &
identical & not tested & identical \\
\texttt{zig\ build-\allowbreak{}exe} (Zig, 1 program) & \textbf{differs (402 bytes)}
& differs (28,099 bytes) & not tested & differs (396 bytes) \\
\texttt{zig\ build-\allowbreak{}exe\ -\allowbreak{}O\ ReleaseSmall} (Zig, 1 program) & identical &
identical & not tested & identical \\
\texttt{ghc\ -O1} (Haskell, 1 program) & identical & identical & not
tested & identical \\
\texttt{cobc\ -x} (COBOL, 1 program) & identical & identical & not
tested & differs (50 bytes) \\
\texttt{ocamlopt\ -\allowbreak{}g} (OCaml, 1 program) & identical & differs (669
bytes) & \textbf{differs (669 bytes)} & identical \\
Julia \texttt{code\_\allowbreak{}llvm} and \texttt{code\_\allowbreak{}native} (1 program) &
\textbf{differs (148 bytes)} & differs (60,594 bytes) & \textbf{differs
(114 bytes)} & differs (100 bytes) \\
\end{longtable}
}

Every cell marked ``identical'' holds for all programs of its
configuration; every cell marked ``differs'' holds for every program of
its configuration. A count is the differing bytes of one program, summed
over its output files where it has several.

\textbf{Findings.}

\begin{enumerate}
\def\labelenumi{\arabic{enumi}.}
\tightlist
\item
  \textbf{Same directory.} On the local machine no configuration
  differed between two builds in the same directory, with the clock
  fixed or free. On the CI runner three did, with the clock fixed:
  \texttt{zig\ build-\allowbreak{}exe} in its default mode, mostly in
  debug-information sections (\texttt{zig\ build-\allowbreak{}exe\ -\allowbreak{}O\ ReleaseSmall}
  was identical); \texttt{cobc\ -x}, in 28 bytes (20 of them in
  \texttt{.\allowbreak{}note.\allowbreak{}gnu.\allowbreak{}build-\allowbreak{}id}); and the Julia run, in 168 bytes of one
  LLVM IR file. For these three configurations the output is not
  determined by the inputs E1 declares. On the macOS laptop seven did:
  every \texttt{-g} configuration of C, C++ and Fortran, in 37 or 38
  bytes of the executable at equal size (in 85 bytes of the executable
  and 16 of the DWARF file for \texttt{gcc\ -\allowbreak{}O2\ -\allowbreak{}g} on
  \texttt{hello.c}); \texttt{zig\ build-\allowbreak{}exe} in its default mode, in 402
  bytes; and the Julia run, in 148 bytes. \texttt{rustc\ -\allowbreak{}O\ -\allowbreak{}g} and the
  other twelve configurations were identical.
\item
  \textbf{Two directories.} On all three machines, every configuration
  that differed between two directories also had the build directory
  written into its output. The configurations identical across
  directories without remapping were \texttt{gcc\ -O2},
  \texttt{rustc\ -O} and \texttt{gfortran\ -\allowbreak{}O2} on all three machines,
  Lean and Wasmtime on the local machine and the macOS laptop,
  \texttt{zig\ build-\allowbreak{}exe\ -\allowbreak{}O\ ReleaseSmall} and \texttt{ghc\ -O1} on the
  CI runner and the macOS laptop, and \texttt{cobc\ -x} on the macOS
  laptop. For the others, the build path is an undeclared input by
  default.
\item
  \textbf{Two directories, remapped.} On the local machine, remapping
  removed every difference except in Go's GOPATH mode, which still
  differed in 58 and 59 bytes, with files of equal size and no copy of
  the build directory left in them. On the CI runner the same
  configuration differed in 59 and 60 bytes, and two more configurations
  differed after remapping: \texttt{gnatmake\ -\allowbreak{}O2\ -\allowbreak{}g}, whose remapped
  outputs still hold two copies of the build directory and differ in its
  two \texttt{.ali} files; and the Julia run, whose remapped run differs
  in 146 bytes of one LLVM IR file with no copy of the build directory
  left. In GOPATH mode with the import path declared, the remapped
  builds were identical across directories on both Linux machines. On
  the macOS laptop, remapping removed no copy of the build directory
  from the \texttt{-g} builds of C, C++, Fortran and Rust: three copies
  are counted in each program's outputs with and without remapping, and
  the differing byte counts with remapping are within one byte of those
  without. Go's GOPATH mode differed after remapping in 119 and 120
  bytes at equal size with no copy of the build directory left,
  \texttt{ocamlopt\ -\allowbreak{}g} in 669 bytes with one copy, and the Julia run in
  114 bytes with none; Go in module mode, and in GOPATH mode with the
  import path declared, was identical across directories when remapped.
\item
  \textbf{Clock.} Only \texttt{gnatmake\ -\allowbreak{}O2\ -\allowbreak{}g} on the CI runner and
  \texttt{cobc\ -x} on the macOS laptop were identical in the same
  directory with \texttt{SOURCE\_\allowbreak{}DATE\_\allowbreak{}EPOCH} set but not with it unset,
  in 3 and 50 bytes; for them the clock is an input.
\item
  \textbf{Across machines.} Between the local machine and the CI runner,
  only configurations whose tool version string is the same on both are
  compared: the gcc, g++, clang, clang++ and gfortran configurations,
  seven configurations on 14 programs, by the SHA-256 of every output of
  one build on each machine, in directories with different absolute
  paths. With remapping, gcc, g++ and gfortran outputs were identical
  across the two machines. \texttt{clang\ -\allowbreak{}O2\ -\allowbreak{}g} and
  \texttt{clang++\ -\allowbreak{}O2\ -\allowbreak{}g} differed across machines for all four of
  their programs even with remapping, although on each machine their
  remapped builds were identical across directories. Go and rustc were
  not compared, because their version strings differ between the two
  machines, and the other toolchains were installed on only one of the
  two. The macOS laptop shares a version string with the CI runner only
  for Zig 0.17.0 and GHC 9.14.1, and with the local machine for no tool.
  The Zig and GHC outputs differ between the CI runner and the laptop,
  which build for different systems and instruction sets (Linux on
  x86-64, macOS on arm64), so that comparison does not test H1 across
  machines.
\end{enumerate}

\textbf{What this means for H1.} The Go GOPATH-mode result is a byte
difference that survives remapping across directories, and the project's
repository records a verdict on it in the registration's table of
changes, dated 2 November 2026 {[}6{]}. Its cause is named. Every
differing byte lies in the \texttt{.\allowbreak{}note.\allowbreak{}go.\allowbreak{}buildid} and
\texttt{.\allowbreak{}note.\allowbreak{}gnu.\allowbreak{}build-\allowbreak{}id} sections, and code and data are identical.
The \texttt{gocachehash} trace shows the compile and link action IDs
hashing the build directory, because GOPATH mode gives a package named
by a file outside every \texttt{GOPATH} the import path \texttt{\_}
followed by its absolute directory, and \texttt{-trimpath} does not
remove it; once the directory is written as a placeholder, the hashed
inputs of the two directories are the same. With only that input
changed, the package placed in a \texttt{GOPATH} and built by its import
path, no hashed input names the directory and the remapped builds are
byte-identical across directories, while the original configuration
still differs. The verdict applies the registered conditions: the
falsifier is a difference that survives remapping when the declared
inputs are equal, and here they were not equal, since the import path is
an input that \texttt{go\ build} derives from the directory. The result
is therefore not a falsifier of H1. The import path can be named and
controlled, so it is added to H1's declared inputs, and H1 holds for Go
in GOPATH mode with the import path named; for a GOPATH-mode build of a
file outside every \texttt{GOPATH}, \texttt{-trimpath} alone is not a
remapping.

The other differences that survive remapping have no recorded verdict.
On the CI runner, \texttt{gnatmake\ -\allowbreak{}O2\ -\allowbreak{}g} and the Julia run differ
across directories after remapping, and E1 does not establish whether a
declared input differs between the two builds. Across machines, the
\texttt{clang\ -\allowbreak{}O2\ -\allowbreak{}g} and \texttt{clang++\ -\allowbreak{}O2\ -\allowbreak{}g} differences after
remapping are differences between builds with the same tool version
string; E1 records no trace of their cause, and the repository has not
recorded a verdict on them. For the macOS laptop, the table of changes
records that the run has no verdict yet, and what it establishes
{[}6{]}. Its remapped builds named the build directory by a path the
tools did not record: the system temporary directory, under
\texttt{/\allowbreak{}var/\allowbreak{}folders}, is reached through the symbolic link
\texttt{/var} to \texttt{/\allowbreak{}private/\allowbreak{}var}, and Go's trace names the
directory with the \texttt{/private} prefix. A check on Linux of gcc,
clang and rustc found that a remapping that names a linked path leaves
the real path in the output. The macOS remapped comparisons therefore do
not test H1's remapping condition. Since that run, E1 names the build
directory by its real path, attributes differing Mach-O bytes to
regions, and adds three configurations that run only on macOS: two in
which clang compiles and links in two calls and the linker is told to
drop the build directory from its debug map, one of them also setting
\texttt{ZERO\_\allowbreak{}AR\_\allowbreak{}DATE=1}, which Apple's linker reads to write zero for
object-file times {[}K{]}; and rustc with the same linker option. No
rerun with these changes is in the repository yet; the table of changes
names that rerun as the evidence for a verdict. Until those causes are
traced, H1 is neither supported nor falsified by E1 as a whole. An
earlier probe in the research sessions, with one program per language,
had found the same-directory and cross-directory pattern for gcc, clang,
rustc, Go in its default mode and gfortran on one machine {[}M{]}.

\textbf{Limits.} The programs are small, one or two per language, so a
difference that only a large program reaches is not tested. Link-time
optimisation, profile-guided builds and parallel code generation beyond
rustc's defaults are not in the matrix. E1 compares bytes; it does not
check that meaning was preserved, which is H3 and H7. The macOS run says
how many bytes differ in each file, not in which part of it.

\subsection{\texorpdfstring{8.2 E2 to E4: protocols, E4's erasure
counts, and E4-Energy (\texttt{reported\_\allowbreak{}j} \& payback
horizons)}{8.2 E2 to E4: protocols, E4's erasure counts, and E4-Energy (reported\_j \& payback horizons)}}\label{e2-to-e4-protocols-e4s-erasure-counts-and-e4-energy-reported_j-payback-horizons}

E2 and E3 have protocols in the repository and have not been run. E4 has
its original 96-row erasure-count run on a host without an energy
counter {[}6{]}, plus a pre-registered Wave 4 Linux RAPL
\texttt{reported\_\allowbreak{}j} extension across 20 deterministic kernels and 14
compiler and runtime tiers (\texttt{board\_\allowbreak{}metered\ =\ false},
\texttt{measured\_\allowbreak{}j\ =\ None}).

\textbf{E2, gate audit (H2).} For each compiler, a fixed set of small
programs, each with one injected fault and a correct twin that differs
only at the fault: a type error for every compiler; a borrow error for
rustc; an unused import or variable for Go; an invalid IR module for
LLVM's \texttt{llc} and an invalid WebAssembly module for
\texttt{wasmtime\ compile}; a failing contract (a Lean theorem whose
proof does not check, an Ada precondition GNATprove cannot prove); and a
generated-code fault for Erlang's validator. Each faulty program and its
twin are built under the default flags and under each flag set
documented as changing the gate (for example \texttt{noEmitOnError} for
tsc, \texttt{fullgraph=True} for TorchDynamo, \texttt{-Werror} for gcc
and clang). The harness records the exit status, every artefact written
with its SHA-256, and the first line of the diagnostic; the twin must
build or the case is void. Each compiler and fault is classified as
\emph{emits nothing by default}, \emph{emits by default and nothing when
configured}, or \emph{emits always}, and each gate's property class,
linking assumption and trusted base is recorded from the compiler's own
documentation, quoted, or as ``not stated''. The falsifier count is the
number of cases where observed and documented behaviour differ. The
tensor compilers' fallbacks to eager execution or to a CPU provider are
recorded as a class of their own. Since the change of 2 November 2026,
each record is evidence of authority pinned (§3.3): it is bound to the
SHA-256 of the program, of the compiler's executable and of every
artefact, and to the flag set, and a record whose digests do not match
its case is void and the case is run again. Each gate's documented
statements also include the authority of the evidence its checker
accepts on and whether that acceptance is recorded bound to digests. A
documented statement is derived evidence and is never promoted to an
observation, so a compiler's class rests on the observed records; a gate
that accepts on inferred evidence, such as a fallback chosen by a
heuristic, is recorded with its acceptance as undecided.

\textbf{E3, equivalence-level audit (H3).} Each pass or flag is assigned
a level before the run: \texttt{-O0}, \texttt{-O2} and \texttt{-O3}
exact; \texttt{-\allowbreak{}ffp-\allowbreak{}contract} settings within a declared error bound;
reassociation (\texttt{-\allowbreak{}fassociative-\allowbreak{}math} with vectorisation,
\texttt{-\allowbreak{}ffast-\allowbreak{}math}, \texttt{-Ofast}) within a declared error bound,
with the bound stated; and, on GPU runners, split reductions and
lowering to FP8 declared lossy. The corpus has integer kernels, exact by
construction, and floating-point kernels with sums, dot products,
polynomial evaluation and reductions, on fixed inputs that include
subnormal numbers, signed zeros, infinities and NaNs (not-a-number
values). The exact reference is the same computation in rational
arithmetic, rounded once at the end; the declared program is the
\texttt{-O0} build with strict IEEE 754 semantics. The metrics are the
bits that differ, the largest ULP distance, the largest absolute and
relative error, and whether a static check (an SMT solver for loop-free
code, or an interval or error-analysis tool) certified the level. The
static tool is still to be chosen. Since the change of 2 November 2026,
each differential-testing record is pinned evidence, bound to the
SHA-256 of the kernel's source, the compiler's executable, the built
binary and the input: an observed change beyond the assigned level is a
counterexample, which is H3's falsifier, and agreement on every input
certifies nothing. A level counts as certified only on proved evidence,
the static tool's verdict bound to the SHA-256 of the binary or IR it
checked, of the tool and of the stated level; a level that only a
heuristic or a model assigns is recorded as undecided. The tensor
compilers and the FP8 and FP4 formats need GPU runners, which the
project does not have yet.

\textbf{E4, erasure and description length against measured energy
(H4).} The corpus is a set of programs, each compiled under several
configurations that E3's checker accepts as equal at the exact level, so
that energy differences come from lowering and not from a change of
program. ε is counted by instrumenting each run: every write of a
register or memory location adds the bits of the previous value that
were not copied out before being overwritten. λ is computed with a proxy
declared before the run. Joules are measured per run, with intervals, in
regime A (an FPGA or in-order core with a supply meter; executed
operations as the baseline), regime B (an x86 host with RAPL, or a GPU
with NVML; time multiplied by power as the baseline) and regime C (a
simulation of entropy production). Compile-time energy is logged too.
The analysis is declared in advance: an ordinary least-squares model per
regime, with the share of variance explained and a 95\% bootstrap
interval over programs; ε and λ explain variance when that interval
excludes zero. Energy figures without a meter reading are labelled
modelled.

\textbf{E4, declared before any energy is read} (2 November 2026)
{[}6{]}. The corpus is eight kernels (\texttt{sumsq}, \texttt{collatz},
\texttt{sieve}, \texttt{sort}, \texttt{crc32}, \texttt{matmul},
\texttt{fib} and \texttt{popcount}), each a function \texttt{run(n)} on
32-bit unsigned integers that uses no library, written in C and in Rust:
16 programs, each built at six optimisation levels (\texttt{clang\ -\allowbreak{}O0},
\texttt{-O1}, \texttt{-O2}, \texttt{-O3}, \texttt{-Os} and \texttt{-Oz};
\texttt{rustc} at \texttt{opt-level} 0, 1, 2, 3, s and z), 96 rows in
all. Because E3's checker does not exist yet, the configurations are
shown equal by a weaker test: each gives the reference value of its
kernel on five declared inputs, natively and as WebAssembly. ε is
counted by an interpreter of WebAssembly written for E4 and pinned by
its SHA-256, on each kernel built for WebAssembly at the row's level; it
is therefore a count for the configuration, not for the native code
whose energy is read, and an upper bound, at byte granularity, on the
bits a run discards. The primary λ is the number of WebAssembly
instructions in the kernel's module. Energy is read in regime B only,
from RAPL on Linux and from \texttt{powermetrics} on macOS; both are
counters reported by the chip or the operating system, not a meter on
the supply, and no host of the harness has a counter for regime A or C.
The analysis uses adjusted R², which can fall below zero, with a 95\%
interval from 2,000 resamples of whole programs; ε and λ explain
variance when that interval lies above zero, and in regime B H4 is
supported only if the interval of the baseline's adjusted R² minus
theirs lies above zero, and falsified if it lies below.

\textbf{E4, counts without energy.} The harness ran on 2 November 2026
(UTC) on a Linux x86-64 machine (Intel Xeon at 2.80 GHz, 4 CPUs, Ubuntu
24.04.5) with clang 18.1.3 and rustc 1.97.0. That machine has no
\texttt{/\allowbreak{}sys/\allowbreak{}class/\allowbreak{}powercap} interface, so every energy field of that
first results file is recorded as unmeasured: the file holds no joules,
no energy analysis is computed from it, and H4 was left untested there.
All 16 programs gave the reference value on every declared input under
every configuration, natively and as WebAssembly, and ε was counted for
all 96 rows. Per call, ε ranges from 272,048 bits (\texttt{popcount} at
\texttt{-O2} and \texttt{-O3} in C, and at \texttt{opt-level} 1, 2 and 3
in Rust) to 10,057,528 bits (Rust \texttt{matmul} at
\texttt{opt-\allowbreak{}level=0}), a modelled Landauer bound at 300 K of 7.81×10⁻¹⁶
J to 2.89×10⁻¹⁴ J per call. Each of the 16 programs has its highest ε at
level 0. Size optimisation does not always lower it: C \texttt{popcount}
at \texttt{-Oz} counts 3,540,016 bits, about 13 times its count at
\texttt{-O2}, and executes 488,404 WebAssembly instructions per call
against 19,530.

\textbf{E4-Energy (Wave 4), Linux RAPL \texttt{reported\_\allowbreak{}j}, compilation
cost and payback horizon \(N^*\).} Before running
(\texttt{PREREG\_\allowbreak{}rce\_\allowbreak{}compile\_\allowbreak{}wave4.\allowbreak{}json}, SHA-256
\texttt{165a12cd\allowbreak{}a4d9c304…}), a 20-kernel deterministic exact-integer and
fixed-point suite was locked across C, Rust, WebAssembly, Python, and
RCE (\texttt{openIE-\allowbreak{}dev/\allowbreak{}rce} at \texttt{1762fb4b}):
\texttt{sha256\_\allowbreak{}block}, \texttt{chacha20\_\allowbreak{}block}, \texttt{q32\_\allowbreak{}orbital},
\texttt{crc32\_\allowbreak{}ieee}, \texttt{lz77\_\allowbreak{}match}, \texttt{dijkstra\_\allowbreak{}grid},
\texttt{ntt\_\allowbreak{}butterfly}, \texttt{matmul\_\allowbreak{}i32}, \texttt{bst\_\allowbreak{}probe},
\texttt{sortnet\_\allowbreak{}16}, \texttt{ringbuf\_\allowbreak{}sim}, \texttt{token\_\allowbreak{}bucket},
\texttt{fnv1a\_\allowbreak{}sip}, \texttt{sieve\_\allowbreak{}primes}, \texttt{collatz\_\allowbreak{}bounded},
\texttt{popcount\_\allowbreak{}swar}, \texttt{edit\_\allowbreak{}levenshtein},
\texttt{rle\_\allowbreak{}codec}, \texttt{modexp\_\allowbreak{}mont} and \texttt{kalman\_\allowbreak{}q16}. All
builds and runs executed on a 32-CPU Linux x86-64 host (AMD Ryzen AI
Max+ 395, Zen 5, Ubuntu Linux 7.0; GCC 15.2.0, Clang 21.1.8, rustc
1.100.0, Node v22.22.1, Python 3.15.0) pinned via
\texttt{sched\_\allowbreak{}setaffinity} to CPU 0, the domain lead CPU of
\texttt{/\allowbreak{}sys/\allowbreak{}class/\allowbreak{}powercap/\allowbreak{}intel-\allowbreak{}rapl:0:0/\allowbreak{}energy\_\allowbreak{}uj} (core
\texttt{MSR\_\allowbreak{}CORE\_\allowbreak{}ENERGY\_\allowbreak{}STAT}), while simultaneously logging
\texttt{/\allowbreak{}sys/\allowbreak{}class/\allowbreak{}powercap/\allowbreak{}intel-\allowbreak{}rapl:0/\allowbreak{}energy\_\allowbreak{}uj} (APU package;
\texttt{basis\ =\ "reported"}, \texttt{board\_\allowbreak{}metered\ =\ false},
\texttt{measured\_\allowbreak{}j\ =\ None}). All 280 kernel-by-tier cells (20 kernels
× 14 tiers) matched their 64-bit unsigned batch checksums over 2,000
seeds with 0 mismatches {[}M{]}.


\begin{center}
\openiefig{fig/rce-compile-fig02-compile-energy-payback.png}{0.94}{https://research.openie.dev/figures/rce-compile/compile-energy-payback.svg}
\end{center}


\emph{Figure 2. Pre-registered Wave 4 E4-Energy and E5-Scale
measurements (Linux RAPL core \texttt{reported\_\allowbreak{}j};
\texttt{board\_\allowbreak{}metered\ =\ false}). (A) Median cumulative lifecycle
energy \(E_{\text{compile}} + N \cdot E_{\text{run}}\) over
\(N \in [1, 10^8]\) kernel invocations. (B) Share of the 20 kernels
where stepping from level 2 (or \texttt{-Os}) to level 3 dominates
immediately (\(N^* = 0\)), pays back at finite \(N^*\), or never
recovers its compile cost (\(\Delta E_{\text{run}} \le 0\)). (C)
E5-Scale soundness across 500 verified semantic mutants and warm De
Bruijn cache savings {[}M{]}.}

{\scriptsize\setlength{\tabcolsep}{2.4pt}\renewcommand{\arraystretch}{1.12}
\begin{longtable}[]{@{}
  >{\raggedright\arraybackslash}p{(\linewidth - 14\tabcolsep) * \real{0.1227}}
  >{\raggedleft\arraybackslash}p{(\linewidth - 14\tabcolsep) * \real{0.1534}}
  >{\raggedleft\arraybackslash}p{(\linewidth - 14\tabcolsep) * \real{0.1534}}
  >{\raggedleft\arraybackslash}p{(\linewidth - 14\tabcolsep) * \real{0.1043}}
  >{\raggedleft\arraybackslash}p{(\linewidth - 14\tabcolsep) * \real{0.0859}}
  >{\raggedleft\arraybackslash}p{(\linewidth - 14\tabcolsep) * \real{0.1472}}
  >{\raggedleft\arraybackslash}p{(\linewidth - 14\tabcolsep) * \real{0.1411}}
  >{\raggedleft\arraybackslash}p{(\linewidth - 14\tabcolsep) * \real{0.0920}}@{}}
\toprule\noalign{}
\begin{minipage}[b]{\linewidth}\raggedright
\textbf{Compiler / Lowering Tier}
\end{minipage} & \begin{minipage}[b]{\linewidth}\raggedleft
\textbf{Median \(E_{\text{compile}}\) Core (\texttt{J})}
\end{minipage} & \begin{minipage}[b]{\linewidth}\raggedleft
\textbf{Median \(E_{\text{compile}}\) Pkg (\texttt{J})}
\end{minipage} & \begin{minipage}[b]{\linewidth}\raggedleft
\textbf{Median Build (\texttt{ms})}
\end{minipage} & \begin{minipage}[b]{\linewidth}\raggedleft
\textbf{Artifact (\texttt{B})}
\end{minipage} & \begin{minipage}[b]{\linewidth}\raggedleft
\textbf{Median \(E_{\text{run}}\) Core (\texttt{µJ})}
\end{minipage} & \begin{minipage}[b]{\linewidth}\raggedleft
\textbf{Median \(E_{\text{run}}\) Pkg (\texttt{µJ})}
\end{minipage} & \begin{minipage}[b]{\linewidth}\raggedleft
\textbf{Median Run (\texttt{ns})}
\end{minipage} \\
\midrule\noalign{}
\endhead
\bottomrule\noalign{}
\endlastfoot
\texttt{python3\ py\_\allowbreak{}compile\ +\ run} & 0.0041 & 0.056 & 0.47 & 3,546 &
241.74 & 2,685.98 & 23,035.4 \\
\texttt{gcc\ -O0} & 0.2904 & 4.892 & 40.16 & 16,384 & 4.17 & 53.24 &
442.5 \\
\texttt{gcc\ -O2} & 0.4723 & 7.528 & 61.51 & 16,320 & 1.81 & 27.24 &
231.5 \\
\texttt{gcc\ -O3} & 0.4955 & 7.890 & 65.78 & 16,320 & 1.75 & 24.36 &
213.7 \\
\texttt{clang\ -\allowbreak{}O0} & 0.4929 & 9.038 & 75.87 & 16,432 & 4.36 & 52.72 &
438.2 \\
\texttt{clang\ -\allowbreak{}Os} & 0.5917 & 10.800 & 88.63 & 16,400 & 1.75 & 22.51 &
190.9 \\
\texttt{clang\ -\allowbreak{}O2} & 0.6391 & 11.459 & 96.09 & 16,368 & 1.27 & 19.61 &
162.3 \\
\texttt{clang\ -\allowbreak{}O3} & 0.7500 & 11.479 & 96.16 & 16,368 & 1.10 & 15.47 &
132.9 \\
\texttt{rustc\ -\allowbreak{}C\ opt-\allowbreak{}level=0} & 0.6450 & 11.639 & 96.58 & 4,567,276 &
13.34 & 159.57 & 1,333.5 \\
\texttt{rustc\ -\allowbreak{}C\ opt-\allowbreak{}level=2} & 1.1040 & 19.838 & 163.65 & 4,533,388 &
1.35 & 19.54 & 171.2 \\
\texttt{rustc\ -\allowbreak{}C\ opt-\allowbreak{}level=3} & 1.1526 & 20.059 & 166.42 & 4,533,552 &
0.79 & 13.29 & 108.5 \\
\texttt{rustc\ wasm32\ -\allowbreak{}O3} (\texttt{node}) & 0.4853 & 8.132 & 67.12 &
1,601 & 1.64 & 20.95 & 157.3 \\
\texttt{rce\ check\ +\ 5-\allowbreak{}target\ emit} (\texttt{node} JS) & 0.8513 &
11.644 & 97.94 & 618,402 & 2,519.99 & 38,343.84 & 319,638.5 \\
\texttt{rce\ check\ +\ emit\ rust\ -\textgreater{}\ rustc\ -O3} &
56.2925 & 896.930 & 7,473.83 & 5,190,208 & 1,306.95 & 15,049.10 &
123,497.6 \\
\end{longtable}
}

\textbf{Optimization payback horizons
\(N^* = \lceil \Delta E_{\text{compile}} / \Delta E_{\text{run}} \rceil\).}
Lowering from Python bytecode to \texttt{clang\ -\allowbreak{}O3} or
\texttt{rustc\ -\allowbreak{}C\ opt-\allowbreak{}level=3} pays back on all 20 kernels after a
median of \textbf{2,490} invocations (range 648 to 10,141) and
\textbf{5,263} invocations (851 to 16,692). Within native compilers,
stepping from unoptimised \texttt{-O0} / \texttt{opt-\allowbreak{}level=0} to
\texttt{-O2} or \texttt{-O3} pays back on all 20 kernels, with median
finite \(N^*\) of \textbf{38,262}
(\texttt{rustc\ 0\ -\textgreater{}\ 2}), \textbf{53,464}
(\texttt{rustc\ 0\ -\textgreater{}\ 3}), \textbf{61,637}
(\texttt{clang\ -O0\ -\textgreater{}\ -O2}), \textbf{81,004}
(\texttt{clang\ -O0\ -\textgreater{}\ -O3}) and \textbf{102,720}
(\texttt{gcc\ -O0\ -\textgreater{}\ -O3}). By contrast, stepping from
\texttt{-O2} (or \texttt{-Os}) to \texttt{-O3} frequently fails to pay
back at any deployment scale: \(\Delta E_{\text{run}} \le 0\)
(\(N^* = \infty\)) on \textbf{6 of 20} kernels (30\%) for
\texttt{clang\ -O2\ -\textgreater{}\ -O3}, \textbf{6 of 20} (30\%) for
\texttt{rustc\ 2\ -\textgreater{}\ 3}, \textbf{7 of 20} (35\%) for
\texttt{clang\ -Os\ -\textgreater{}\ -O3}, and \textbf{4 of 20} (20\%)
for \texttt{gcc\ -O2\ -\textgreater{}\ -O3}, while the kernels that do
gain require a median of \textbf{106,920} to \textbf{244,165}
invocations (and up to \textbf{3.87×10⁶} on Clang \texttt{sha256\_\allowbreak{}block}
and \textbf{1.39×10⁷} on GCC) before \texttt{-O3} recovers its extra
compile joules.
\texttt{rustc\ -\allowbreak{}\/\allowbreak{}-\allowbreak{}target\ wasm32-\allowbreak{}unknown-\allowbreak{}unknown\ -\allowbreak{}C\ opt-\allowbreak{}level=3}
builds a 1.6 KB portable module in 0.485 J core and executes in Node at
1.64 µJ per call, so native \texttt{rustc\ -\allowbreak{}O3} needs a median of
\textbf{871,412} invocations to amortise its extra 0.667 J of native
link and codegen work. For RCE (\texttt{1762fb4b}), \texttt{rce\ check}
plus emitting all five scalar targets (\texttt{core}, \texttt{rust},
\texttt{python}, \texttt{js}, \texttt{scala}) costs a median of 0.851 J
core (97.9 ms), whereas compiling the standalone 1,400-line RCE
value-runtime prelude inside each emitted Rust file under
\texttt{rustc\ -\allowbreak{}O3} costs 56.29 J core and requires a median of
\textbf{31,080} invocations (18 of 20 finite, 2 never) to beat running
the emitted JS directly in Node.

\textbf{H4 Regime-B check on Linux RAPL.} Across the 140 Clang and Rust
rows with freestanding WebAssembly IR modules (40 kernel-by-language
series, 2,000 whole-kernel bootstrap resamples), static WebAssembly IR
description length \(\lambda\) (\texttt{ir\_nodes} and
\texttt{wasm\_\allowbreak{}bytes}) explains an adjusted \(R^2\) of \textbf{0.1154}
(95\% interval \textbf{{[}0.0654, 0.4058{]}}), while the Regime-B
baseline (\texttt{run\_\allowbreak{}wall\_\allowbreak{}s\ ×\ mean\_\allowbreak{}core\_\allowbreak{}power\_\allowbreak{}w}) explains
\textbf{0.9999} (\textbf{{[}0.9987, 1.0000{]}}), a positive gap of
\textbf{+0.8845} (\textbf{{[}+0.5979, +0.9352{]}}) {[}M{]}. That matches
H4's Regime-B prediction on Linux RAPL, alongside the earlier macOS
\texttt{macmon} runs (\texttt{HYPOTHESES.\allowbreak{}md}); H4 as a whole remains
open until Regimes A and C are tested and quiet dedicated hosts are
compared against shared-host runs (non-pinned cores on this host carried
co-tenant load averages of 18.97 to 36.42).

\textbf{E6 and E7} (H6 and H7) are planned, with no protocol yet. E6
needs a compiler built on the gear ladder, which does not exist. E5 has
run in two stages, and §8.3 reports both.

\subsection{8.3 E5: a name-free hash,
measured}\label{e5-a-name-free-hash-measured}

\textbf{Method.} E5 tests H5 and implements part R of the Compiler Form
as two operations {[}6{]}. Its protocol and harness are dated 2 November
2026, before its first run, and the registration records E5's
predictions and falsifiers under H5 on the same date. Each operation
hashes a build output after removing what names the identifiers a
program binds.

\begin{itemize}
\tightlist
\item
  \emph{At the machine-code level}, a function identifier in the style
  of WARP, following the parts of WARP's published description that the
  E5 protocol quotes, is computed over unlinked object files: x86-64 ELF
  objects, arm64 Mach-O objects and WebAssembly objects. Each function
  is split into basic blocks; every instruction that a relocation record
  (a note telling the linker to patch an address there) points into is
  replaced by zero bytes; no-operation instructions and moves of a
  full-width register to itself are left out; each block is hashed, and
  the function's identifier is a hash of its blocks' hashes. Two keys
  come from it: the identifier of the kernel's function \texttt{run},
  and a hash of the sorted identifiers of every function in the object,
  which no name enters.
\item
  \emph{At the IR level}, a hash of linked WebAssembly modules modulo
  alpha-equivalence (two programs are alpha-equivalent when a one-to-one
  renaming of bound identifiers turns one into the other) writes each
  module in a canonical form and hashes it with SHA-256. In that form a
  type index becomes the signature it names, locals are numbered by
  position and first use, functions are ordered by the shape of the call
  graph, independent definitions are sorted, and custom sections, such
  as the one that holds names, are left out. One key covers the whole
  module and one the part that \texttt{run} reaches. The protocol
  credits the idea of hashing modulo alpha-equivalence to Maziarz and
  colleagues (PLDI 2021) and states that its canonical form is not their
  algorithm.
\end{itemize}

Two baseline keys stand for the two kinds of cache key described under U
in §3.5: the SHA-256 of the output bytes, and the SHA-256 of the source
text and build command with the build directory removed. H5's metric
names the compiler's own key; E5 reads no compiler's cache key and uses
the two baselines in its place, a substitution the table of changes
records. E5 compares keys; it runs no cache.

\textbf{Corpus and pairs.} The corpus is E4's 16 programs at their six
levels, each with four outputs (an x86-64 object, an arm64 Mach-O object
for C only, a WebAssembly object, and E4's linked WebAssembly module),
and E1's C and Rust programs (\texttt{hello.c}, \texttt{sort.c},
\texttt{hello.rs} and \texttt{generic.\allowbreak{}rs}) at the same six levels, as
x86-64 objects only. Each source is built from directory A, its same
bytes from directory B, and a renamed copy from directory N; the three
paths have equal length. The renamer gives every identifier the program
binds a new name derived from the SHA-256 of the old one, and removes
comments. The evidence that a renamed source is the same program is its
behaviour: the linked modules of all three variants give E4's reference
value on every declared input (96 of 96 in each variant), and each of
E1's four programs prints the same bytes before and after renaming. A
against B (\emph{two directories}) and A against N (\emph{renamed}) are
declared equal. Every pair of levels of one toolchain, the C and Rust
builds of one kernel at matching levels, and every pair of builds of two
different kernels are declared not equal, as are the \emph{mutations}:
each linked module against itself with one opcode changed to another of
its class, one constant changed or one memory offset changed.

A \emph{hit} is an equal key on a two-directory or renamed pair. A
\emph{false hit} is an equal key on builds of two different kernels, or
an IR hash left unchanged by a mutation. An equal key on two levels, or
on the C and Rust builds of one kernel, is reported and not counted as
false, because each such pair computes one function. H5 is supported if,
on the renamed pairs, every R hash has a higher hit rate than both
baselines and there is no false hit; it is falsified by any false hit,
or if no R hash has a higher hit rate than both baselines on the renamed
pairs.

\textbf{Machine.} E5 ran on 2 November 2026 (UTC) on a Linux x86-64
machine (Intel Xeon at 2.80 GHz, 4 CPUs, Ubuntu 24.04.5) with clang
18.1.3 and rustc 1.97.0. It built and hashed 360 x86-64 objects, 144
Mach-O objects, 288 WebAssembly objects and 288 linked modules. The
instruction boundaries found by its x86-64 decoder agreed with those of
LLVM's \texttt{llvm-\allowbreak{}objdump} in all 386 functions of the x86-64 objects
built in directory A.

\textbf{Results.} Equal keys, out of the pairs compared. The R hash is
the machine-code identifier for objects and the IR hash for linked
modules.

{\scriptsize\setlength{\tabcolsep}{3.0pt}\renewcommand{\arraystretch}{1.14}
\begin{longtable}[]{@{}
  >{\raggedright\arraybackslash}p{(\linewidth - 12\tabcolsep) * \real{0.1471}}
  >{\raggedright\arraybackslash}p{(\linewidth - 12\tabcolsep) * \real{0.2157}}
  >{\raggedright\arraybackslash}p{(\linewidth - 12\tabcolsep) * \real{0.0784}}
  >{\raggedright\arraybackslash}p{(\linewidth - 12\tabcolsep) * \real{0.1373}}
  >{\raggedright\arraybackslash}p{(\linewidth - 12\tabcolsep) * \real{0.1176}}
  >{\raggedright\arraybackslash}p{(\linewidth - 12\tabcolsep) * \real{0.1765}}
  >{\raggedright\arraybackslash}p{(\linewidth - 12\tabcolsep) * \real{0.1275}}@{}}
\toprule\noalign{}
\begin{minipage}[b]{\linewidth}\raggedright
\textbf{Comparison}
\end{minipage} & \begin{minipage}[b]{\linewidth}\raggedright
\textbf{Output}
\end{minipage} & \begin{minipage}[b]{\linewidth}\raggedright
\textbf{Pairs}
\end{minipage} & \begin{minipage}[b]{\linewidth}\raggedright
\textbf{Output bytes}
\end{minipage} & \begin{minipage}[b]{\linewidth}\raggedright
\textbf{Source key}
\end{minipage} & \begin{minipage}[b]{\linewidth}\raggedright
\textbf{R hash, whole output}
\end{minipage} & \begin{minipage}[b]{\linewidth}\raggedright
\textbf{R hash, \texttt{run}}
\end{minipage} \\
\midrule\noalign{}
\endhead
\bottomrule\noalign{}
\endlastfoot
Two directories & x86-64 objects, E4's kernels & 96 & 87 & 96 & 96 &
96 \\
Two directories & x86-64 objects, E1's programs & 24 & 24 & 24 & 24 & no
\texttt{run} \\
Two directories & arm64 Mach-O objects & 48 & 48 & 48 & 48 & 48 \\
Two directories & WebAssembly objects & 96 & 87 & 96 & 96 & 96 \\
Two directories & linked WebAssembly modules & 96 & 87 & 96 &
\textbf{87} & \textbf{87} \\
Renamed & x86-64 objects, E4's kernels & 96 & 44 & 0 & 96 & 96 \\
Renamed & x86-64 objects, E1's programs & 24 & 17 & 12 & 24 & no
\texttt{run} \\
Renamed & arm64 Mach-O objects & 48 & 24 & 0 & 48 & 48 \\
Renamed & WebAssembly objects & 96 & 44 & 0 & 96 & 96 \\
Renamed & linked WebAssembly modules & 96 & 76 & 0 & \textbf{87} &
\textbf{87} \\
\end{longtable}
}

\textbf{Findings.}

\begin{enumerate}
\def\labelenumi{\arabic{enumi}.}
\tightlist
\item
  \textbf{No false hit.} No R hash was equal on any of 26,208
  comparisons of builds of two different kernels (4,032 pairs per output
  and key, 1,008 for the Mach-O objects), and none of 9,328 mutations of
  96 linked modules (4,712 of an opcode, 3,564 of a constant and 1,052
  of a memory offset) left the IR hash unchanged.
\item
  \textbf{A gain on renamed pairs.} On the renamed pairs every R hash,
  in each of its nine combinations of output and key, had a higher hit
  rate than both baselines: 100\% for the machine-code identifiers,
  against 45.8\% to 70.8\% for output bytes and 0\% to 50.0\% for the
  source key; 90.6\% for the IR hashes, against 79.2\% and 0\%.
\item
  \textbf{The verdict.} By its registered conditions, H5 is supported on
  this corpus. The table of changes records the verdict together with
  the predictions that did not hold, and limits it to this corpus, two
  toolchains and one IR {[}6{]}.
\item
  \textbf{Misses at the IR level.} The IR hashes were unequal on 9 of
  the 96 two-directory pairs and 9 of the 96 renamed pairs: the linked
  Rust modules of all eight kernels at \texttt{opt-\allowbreak{}level=0}, and Rust
  \texttt{sort} at \texttt{opt-\allowbreak{}level=1}. In each, only the data part of
  the canonical form differs, and each module holds one copy of its
  build directory's path. The registered IR hash treats data bytes as
  constants and does not remove the build path, which the Form's R
  removes. A diagnostic key added after the first run, which writes that
  path in the data as a placeholder and is not used for the verdict, is
  equal on all 96 two-directory pairs and on 92 of the 96 renamed pairs.
  The four left are Rust \texttt{fib}, \texttt{matmul}, \texttt{sieve}
  and \texttt{sort} at \texttt{opt-\allowbreak{}level=0}, again in their data. For
  \texttt{fib}, a check outside the harness found the two data segments
  differing in the column numbers of the location records that Rust's
  overflow messages carry: the renamed identifiers are longer, so the
  same expression starts at a later column.
\item
  \textbf{Other predictions that did not hold.} The source key was equal
  on 12 of the 24 renamed pairs of E1's programs, because the two hello
  programs bind no identifier and the renamer left them unchanged. The
  machine-code identifier of \texttt{run} in the WebAssembly objects was
  equal for the C and the Rust build of \texttt{fib} at five matching
  levels (1, 2, 3, s and z) and of \texttt{popcount} at s, 6 of 48
  pairs; each such pair computes one function, so it is not a false hit.
  On the linked modules the IR hashes differ in all 48 C and Rust pairs;
  the functions part of their canonical forms is equal in 6 of the 48,
  and the memory part in none.
\item
  \textbf{Levels.} Across pairs of optimisation levels of one toolchain,
  the R hashes were equal on 10.0\% to 21.7\% of pairs and the output
  bytes on 10.0\% to 12.5\%, as registered: unequal on most pairs, with
  equal pairs reported and not counted.
\end{enumerate}

\textbf{What E5 does not show.} An equal machine-code identifier
establishes no equivalence on its own: a zeroed instruction can call or
load a different symbol once the object is linked, and the identifier
does not see which. WARP's constraints, which tell such functions apart,
are not implemented, so a machine-code hit is a candidate for reuse that
a checker still has to accept, and the protocol gives these identifiers
no property class. An equal IR hash means that the two linked modules
are alpha-equivalent, barring a SHA-256 collision, which the protocol
takes as the exact level for their integer results under the WebAssembly
semantics; it says nothing about the native code built from the same
source. Unequal hashes say nothing about inequivalence: pairs of levels
compute one function under different hashes, and two alpha-equivalent
modules whose functions tie in the canonical order could be ordered
differently, since ties are counted, not resolved. If an IR hash keyed
reuse at rung G0, the protocol's gate statement would be equivalence at
the exact level of the module, linked by its import and export names by
a host that does not read custom sections, with E5's own readers,
canonicaliser and hashing code, Python's hashing libraries and the
WebAssembly specification as the trusted base. That no false hit
occurred among these pairs and mutations does not show that the hashes
have none on other programs: the initial E5 corpus is 16 kernels and 4
programs, from two toolchains that both lower through LLVM, run on one
machine, and the IR is WebAssembly, not LLVM IR. And since E5 reads no
compiler's own cache key, it does not show that any compiler's cache
would gain hits.

\textbf{E5-Scale (Wave 4): 185 \texttt{W1} RCE Forms, 60 C/Rust/Python
kernels, 500 semantic mutants, and avoided \texttt{reported\_\allowbreak{}j}.} To
test H5 beyond 20 small files and measure the compilation joules a
name-free cache avoids, E5 was scaled by more than 12×
(\texttt{PREREG\_\allowbreak{}rce\_\allowbreak{}compile\_\allowbreak{}wave4.\allowbreak{}json}) across \textbf{245 base
programs}: all \textbf{185} emitted Forms of the \texttt{W1} RCE Form
corpus
(\texttt{openie-\allowbreak{}research-\allowbreak{}experiments/\allowbreak{}rce/\allowbreak{}bench/\allowbreak{}form-\allowbreak{}space/\allowbreak{}results/\allowbreak{}forms.\allowbreak{}jsonl},
evaluated with \texttt{rce\ fmt\ nameless} at \texttt{1762fb4b}) plus
the \textbf{60} C, Rust and Python kernel sources of E4-Energy
(\texttt{20\ ×\ 3}). Four cache keys were compared across
formatting/comment shifts, systematic alpha-renaming, benign
independent-constant or event reordering, \textbf{29,890} unordered
cross-program pairs, and \textbf{500} single-site semantic mutations
(150 variable-binding swaps such as \texttt{x\ -\ y} to
\texttt{y\ -\ x}, 150 binary-operator changes, 125 constant changes, and
75 branch/guard changes), every one verified non-equivalent on concrete
execution vectors before hashing {[}M{]}.

{\scriptsize\setlength{\tabcolsep}{3.0pt}\renewcommand{\arraystretch}{1.14}
\begin{longtable}[]{@{}
  >{\raggedright\arraybackslash}p{(\linewidth - 10\tabcolsep) * \real{0.2773}}
  >{\raggedleft\arraybackslash}p{(\linewidth - 10\tabcolsep) * \real{0.1261}}
  >{\raggedleft\arraybackslash}p{(\linewidth - 10\tabcolsep) * \real{0.1092}}
  >{\raggedleft\arraybackslash}p{(\linewidth - 10\tabcolsep) * \real{0.1849}}
  >{\raggedleft\arraybackslash}p{(\linewidth - 10\tabcolsep) * \real{0.1513}}
  >{\raggedleft\arraybackslash}p{(\linewidth - 10\tabcolsep) * \real{0.1513}}@{}}
\toprule\noalign{}
\begin{minipage}[b]{\linewidth}\raggedright
\textbf{Corpus / Adversarial Set}
\end{minipage} & \begin{minipage}[b]{\linewidth}\raggedleft
\textbf{Pairs / Mutants}
\end{minipage} & \begin{minipage}[b]{\linewidth}\raggedleft
\textbf{Raw SHA-256}
\end{minipage} & \begin{minipage}[b]{\linewidth}\raggedleft
\textbf{Comment/Whitespace-Stripped}
\end{minipage} & \begin{minipage}[b]{\linewidth}\raggedleft
\textbf{Naive AST ID-Strip}
\end{minipage} & \begin{minipage}[b]{\linewidth}\raggedleft
\textbf{De Bruijn / Nameless}
\end{minipage} \\
\midrule\noalign{}
\endhead
\bottomrule\noalign{}
\endlastfoot
\texttt{185} \texttt{W1} RCE Forms · alpha-rename
(\texttt{rce\ fmt\ nameless}) & 185 hits & 0 (0.0\%) & 50 (27.0\%) & 185
(100.0\%) & \textbf{128 (69.2\%)} \\
\texttt{60} C/Rust/Python kernels · formatting \& comment shift & 60
hits & 0 (0.0\%) & 60 (100.0\%) & 60 (100.0\%) & \textbf{60
(100.0\%)} \\
\texttt{60} C/Rust/Python kernels · alpha-rename & 60 hits & 0 (0.0\%) &
0 (0.0\%) & 60 (100.0\%) & \textbf{40 (66.7\%)} \\
\texttt{60} C/Rust/Python kernels · constant-order swap & 60 hits & 0
(0.0\%) & 0 (0.0\%) & 0 (0.0\%) & \textbf{60 (100.0\%)} \\
\textbf{Soundness:} \texttt{150} variable-binding swap mutants (false
hits) & 150 bugs & 0 & 0 & \textbf{150 (100.0\%, unsound)} & \textbf{0
(0.00\%)} \\
\textbf{Soundness:} \texttt{350} operator + literal + guard mutants
(false hits) & 350 bugs & 0 & 0 & 0 (0.0\%) & \textbf{0 (0.00\%)} \\
\textbf{Soundness:} All \texttt{500} verified semantic mutants (false
hits) & 500 bugs & 0 (0.0\%) & 0 (0.0\%) & \textbf{150 (30.0\%,
unsound)} & \textbf{0 (0.00\%)} \\
\textbf{Soundness:} All unordered cross-program pairs (\texttt{245}
bases) & 29,890 pairs & 0 & 0 & --- & \textbf{0 (0.00\%)} \\
\end{longtable}
}

Three empirical results stand out:

\begin{enumerate}
\def\labelenumi{\arabic{enumi}.}
\tightlist
\item
  \textbf{Why De Bruijn positions beat naive identifier-stripping.}
  Replacing every user identifier with a single anonymous \texttt{ID}
  token gives a 100\% hit rate on alpha-renaming, but \textbf{collides
  on 150 of 150 (100.0\%) variable-binding swap bugs} (\texttt{x\ -\ y}
  and \texttt{y\ -\ x} both project to \texttt{ID\ -\ ID}), yielding a
  30.0\% overall false-hit rate across the 500 verified bugs. Replacing
  each bound occurrence by its binding slot (\texttt{v0\ -\ v1} vs
  \texttt{v1\ -\ v0}) yields \textbf{0 of 500 (0.00\%)} false-positive
  collisions and \textbf{0 of 29,890} cross-program collisions,
  supporting H5 at scale.
\item
  \textbf{Real formatter behaviour of \texttt{rce\ fmt\ nameless} at
  \texttt{1762fb4b}.} Because \texttt{rce\ fmt\ nameless} is designed to
  emit valid, checkable RCE source text with comments preserved
  (\texttt{nameless.\allowbreak{}rs}), its raw stdout changes when \texttt{\#}
  comments are inserted (0 of 185 on comment-added formatting shifts,
  and 128 of 185 on alpha-renaming where 50 zero-binder or foreign-field
  Forms fell back to comment shifts) and when adjacent \texttt{event}
  declarations swap order (0 of 23, since \texttt{e0,\ e1} follow
  declaration order); stripping \texttt{\#} comments before hashing
  \texttt{rce\ fmt\ nameless} restores 185 of 185 (100\%) on comment
  shifts.
\item
  \textbf{Avoided compilation \texttt{reported\_\allowbreak{}j}.} Sitting a warm De
  Bruijn source cache in front of alpha-renamed rebuilds on pinned CPU 0
  lowers per-kernel work from \textbf{1.166 J} core (143.8 ms) to
  \textbf{2.60 mJ} (288.1 µs) on \texttt{rustc\ -\allowbreak{}C\ opt-\allowbreak{}level=3}
  (\textbf{448.5×} reduction, avoiding 1.163 J core per hit), from
  \textbf{0.682 J} core (82.0 ms) to \textbf{4.07 mJ} (342.0 µs) on
  \texttt{clang\ -\allowbreak{}O3} (\textbf{167.8×} reduction, avoiding 0.678 J core
  per hit), and from \textbf{78.71 mJ} core (9.85 ms) to \textbf{4.59
  mJ} (0.74 ms) for \texttt{rce\ check} plus 5-target \texttt{rce\ emit}
  (\textbf{17.15×} reduction) {[}M{]}. Readers can test live snippets
  and inspect all 280 E4-Energy cells in the companion observatory at
  \href{https://research.openie.dev/living/rce-compile/}{\texttt{/\allowbreak{}living/\allowbreak{}rce-\allowbreak{}compile/\allowbreak{}}}.
\end{enumerate}

\section{9. Limits of this study}\label{limits-of-this-study}

\begin{itemize}
\tightlist
\item
  \textbf{Almost all external evidence is {[}S{]} or {[}K{]}.} The hosts
  that hold the primary texts could not be reached from the research
  sessions, the GitHub commits API with a path filter was refused, and
  the national databases (ACM DL, CNKI, J-STAGE, eLibrary.ru, HAL) were
  reached only through general web search. Before publication beyond a
  draft, the primary texts need reading, in this order: Hoare, He and
  Sampaio (1993); Ershov (1977); the theorem statement in the CompCert
  manual; arXiv:2610.33193, to check PP-complete against PP-hard;
  arXiv:2609.11356; and the Zig release notes.
\item
  \textbf{The negative search for a named law is open in some
  languages.} Native terms for ``law of compilation'' in Chinese,
  Russian, French and German were not searched one by one.
\item
  \textbf{The per-part table is qualitative.} It summarises architecture
  and documentation; it is not a measurement of each compiler.
\item
  \textbf{Energy readings are Regime-B on-chip \texttt{reported\_\allowbreak{}j} on a
  shared host, not supply wattmeter \texttt{measured\_\allowbreak{}j}.} E4-Energy
  records Linux RAPL core and package \texttt{reported\_\allowbreak{}j}
  (\texttt{board\_\allowbreak{}metered\ =\ false}, \texttt{measured\_\allowbreak{}j\ =\ None}) on
  a single 32-CPU Zen 5 host whose non-pinned cores carried co-tenant
  workloads (§8.2); Regimes A (in-order/FPGA supply shunt) and C
  (thermodynamic entropy production) remain unmeasured, so H4 as a
  three-regime hypothesis is not yet closed.
\item
  \textbf{Corpus scope.} E1 builds one or two programs per language;
  E4-Energy covers 20 self-contained integer and fixed-point kernels
  without multi-crate link-time optimisation or GPU kernels; and
  E5-Scale's source-level De Bruijn canonicaliser covers RCE Forms
  (\texttt{185}) plus single-file C, Rust and Python kernels
  (\texttt{60}), not full multi-translation-unit C++ or LLVM IR (§8.3).
\item
  \textbf{The 98/2 split is registered, not yet measured.} It is
  hypothesis H6 (§5, §7), and its measurement is planned rung by rung
  with a declared unit.
\end{itemize}

\section{10. Conclusion}\label{conclusion}

Measured against real compilers, the candidate splits into two kinds of
part. Representation, cost, checker, reuse and gate describe practice
that exists, unevenly, somewhere in the field. Lowering, the convergence
function and energy describe a standard no shipping compiler meets. That
split is the reason to restate the candidate as a normal form: placing
gcc, rustc, Go, GHC, XLA and the rest against one declared contract,
with path-stripped determinism, named gates, explicit equivalence levels
and a measured distance from two computable bounds, gives compilers from
these communities a common scale. Two empirical gaps in the initial
draft now carry Regime-B data. First, E5-Scale shows across 245 programs
and 500 verified semantic mutants why part R must replace bound
identifiers by De Bruijn positions rather than merely stripping names:
positional hashing yields 0 of 500 false hits and cuts cold
\texttt{clang\ -\allowbreak{}O3} and \texttt{rustc\ -\allowbreak{}O3} rebuild \texttt{reported\_\allowbreak{}j}
by 167.8× and 448.5×, whereas naive identifier-stripping collides on
100\% of variable-binding swaps (§8.3). Second, E4-Energy shows across
20 exact-integer kernels and 14 compiler tiers that compilation joules
materially reshape which optimisation tier wins: lowering from Python to
\texttt{clang\ -\allowbreak{}O3} pays back after a median of 2,490 runs, \texttt{-O0}
to \texttt{-O3} needs 53,464 to 102,720 runs, \texttt{-O2} to
\texttt{-O3} never recovers its extra compile joules on 20\% to 30\% of
the kernels, and in Regime B time multiplied by power explains far more
variance in RAPL \texttt{reported\_\allowbreak{}j} than static WebAssembly IR size
(0.9999 against 0.1154, §8.2). With hardware \(10^6\) to \(10^{11}\)
above the floor, ranking compilers by lifecycle \texttt{reported\_\allowbreak{}j}
(and supply \texttt{measured\_\allowbreak{}j} when hardware shunt meters are
attached) while keeping Landauer as a physical lower bound turns part E
into an operational decision rule. The 98/2 principle needs the same
treatment: the unit it is counted in must be chosen and published before
the first measurement.

\section{Appendix A. Claims checked}\label{appendix-a.-claims-checked}

The discussion in which the candidate was proposed (§1) made 55 factual
claims about the field. Each was checked against the sources the
research sessions could reach. ``Confirmed'' means the source says what
the claim says; ``partly'' means the source exists but the claim
overstates, misdates or misattributes it; ``contradicted'' means the
source says something incompatible; ``not found'' means no source was
located. Nearly every confirmation is {[}S{]}, so ``confirmed'' means
confirmed against snippet text.

{\scriptsize\setlength{\tabcolsep}{3.0pt}\renewcommand{\arraystretch}{1.14}
\begin{longtable}[]{@{}
  >{\raggedright\arraybackslash}p{(\linewidth - 10\tabcolsep) * \real{0.3299}}
  >{\raggedright\arraybackslash}p{(\linewidth - 10\tabcolsep) * \real{0.1134}}
  >{\raggedright\arraybackslash}p{(\linewidth - 10\tabcolsep) * \real{0.1340}}
  >{\raggedright\arraybackslash}p{(\linewidth - 10\tabcolsep) * \real{0.1134}}
  >{\raggedright\arraybackslash}p{(\linewidth - 10\tabcolsep) * \real{0.1753}}
  >{\raggedright\arraybackslash}p{(\linewidth - 10\tabcolsep) * \real{0.1340}}@{}}
\toprule\noalign{}
\begin{minipage}[b]{\linewidth}\raggedright
\textbf{Topic}
\end{minipage} & \begin{minipage}[b]{\linewidth}\raggedright
\textbf{Claims}
\end{minipage} & \begin{minipage}[b]{\linewidth}\raggedright
\textbf{Confirmed}
\end{minipage} & \begin{minipage}[b]{\linewidth}\raggedright
\textbf{Partly}
\end{minipage} & \begin{minipage}[b]{\linewidth}\raggedright
\textbf{Contradicted}
\end{minipage} & \begin{minipage}[b]{\linewidth}\raggedright
\textbf{Not found}
\end{minipage} \\
\midrule\noalign{}
\endhead
\bottomrule\noalign{}
\endlastfoot
Correctness and verification & 9 & 4 & 3 & 0 & 2 \\
Optimisation and model-assisted compilation & 11 & 2 & 3 & 2 & 4 \\
Physics and thermodynamic hardware & 9 & 4 & 3 & 0 & 2 \\
Systems languages and industry & 10 & 4 & 3 & 1 & 2 \\
HPC, legacy and high-integrity compilers & 7 & 4 & 0 & 0 & 3 \\
ML compilers & 8 & 2 & 2 & 1 & 3 \\
The law itself & 1 & 0 & 0 & 0 & 1 \\
\textbf{Total} & \textbf{55} & \textbf{20} & \textbf{14} & \textbf{4} &
\textbf{17} \\
\end{longtable}
}

\textbf{Figures were wrong more often than existence.} Most named
papers, products and deals exist, but the numbers attached to them were
misattributed, rounded in the claim's favour, moved between comparisons,
or not found at all. The thermodynamic-compilation thread is the
weakest: its headline paper and its largest numbers were not found, and
the nearest matches are self-published preprints.

\subsection{A.1 Contradicted claims, with the
correction}\label{a.1-contradicted-claims-with-the-correction}

{\footnotesize\setlength{\tabcolsep}{3.8pt}\renewcommand{\arraystretch}{1.15}
\begin{longtable}[]{@{}
  >{\raggedright\arraybackslash}p{(\linewidth - 4\tabcolsep) * \real{0.3918}}
  >{\raggedright\arraybackslash}p{(\linewidth - 4\tabcolsep) * \real{0.5361}}
  >{\raggedright\arraybackslash}p{(\linewidth - 4\tabcolsep) * \real{0.0722}}@{}}
\toprule\noalign{}
\begin{minipage}[b]{\linewidth}\raggedright
\textbf{Claim}
\end{minipage} & \begin{minipage}[b]{\linewidth}\raggedright
\textbf{Correction}
\end{minipage} & \begin{minipage}[b]{\linewidth}\raggedright
\textbf{Mark}
\end{minipage} \\
\midrule\noalign{}
\endhead
\bottomrule\noalign{}
\endlastfoot
TENSAT gives XLA 3.45\% on average and 56\% on NasRNN & The figures come
from a different paper, ``Mind the Abstraction Gap'' (OOPSLA 2025)
{[}19{]}; TENSAT reports up to 16\% {[}18{]} & {[}S{]} \\
Japan's ``SlimeCompiler'' uses non-commutative rings for a 1.5× to 3×
gain & A self-published preprint (December 2025) with a benchmark
labelled ``simulated'' and one program at 1.2×; the method is
read/write-set dependence analysis; no Japanese origin was found
{[}56{]} & {[}S{]} \\
Mojo went fully open source on 23 October 2026 & It was open-sourced on
18 August 2026; whether the MAX platform is covered is disputed {[}57{]}
& {[}S{]} \\
Triton has 300 to 500 commits a month & 159 to 267 a month, April to
October 2026 {[}58{]} & {[}F{]} \\
\end{longtable}
}

\subsection{A.2 Claims that partly hold, with the
correction}\label{a.2-claims-that-partly-hold-with-the-correction}

{\footnotesize\setlength{\tabcolsep}{3.8pt}\renewcommand{\arraystretch}{1.15}
\begin{longtable}[]{@{}
  >{\raggedright\arraybackslash}p{(\linewidth - 4\tabcolsep) * \real{0.3874}}
  >{\raggedright\arraybackslash}p{(\linewidth - 4\tabcolsep) * \real{0.4685}}
  >{\raggedright\arraybackslash}p{(\linewidth - 4\tabcolsep) * \real{0.1441}}@{}}
\toprule\noalign{}
\begin{minipage}[b]{\linewidth}\raggedright
\textbf{Claim}
\end{minipage} & \begin{minipage}[b]{\linewidth}\raggedright
\textbf{Correction}
\end{minipage} & \begin{minipage}[b]{\linewidth}\raggedright
\textbf{Mark}
\end{minipage} \\
\midrule\noalign{}
\endhead
\bottomrule\noalign{}
\endlastfoot
Bryant's 2024 UBC thesis shows that compositional correctness reduces to
confluence & The thesis recasts correctness of one ANF compiler model as
confluence and proves ``several related but weaker properties'' {[}14{]}
& {[}S{]} \\
October 2026 work on weakest-precondition-preserving compilation,
mechanised in Rocq & A PriSC 2026 talk at POPL in January 2026 {[}59{]};
the Rocq mechanisation and the October date were not found & {[}S{]},
{[}NOT FOUND{]} \\
``The Next 700 Compiler Correctness Theorems'' unifies compiler
correctness & It gives a framework for stating and comparing theorems
and argues that ``correct'' is ambiguous until the linking context is
named; it is not one theorem {[}10{]} & {[}S{]} \\
Meta's coreset pass ordering beats -Oz by about 4.5\% on about 5,000
repositories & 4.7\% on 4,683 repositories {[}60{]} & {[}S{]} \\
``Compiled AI'' (Stanford, April 2026): 96\%, break-even at 17 runs, a
57-fold saving at 1,000 runs & Led by XY.AI Labs; 96\% and about 17 runs
confirmed; the 57-fold figure was not found {[}42{]} & {[}S{]}, {[}NOT
FOUND{]} \\
Model-written Rocq proofs of crash-freedom (October 2026) at about \$40
per kilobyte, with 30-fold reuse savings & About 180 million lines of
Rocq for 239 OSS-Fuzz entry points in 123 libraries, for inputs under 64
MB {[}61{]}; the cost and reuse figures were not found, and Anthropic's
research index lists no such post {[}62{]} & {[}S{]}, {[}F{]} \\
Berkeley Lab's generative thermodynamic computing gives an advantage
above \(10^{11}\) & The paper exists and is shown in digital simulation
{[}63{]}; the \(10^{11}\) figure was not found & {[}S{]}, {[}NOT
FOUND{]} \\
Normal Computing's CN101 thermodynamic chip & Taped out July 2025, shown
at Hot Chips 2026; it is digital pseudorandom CMOS, and no measured
joules were found {[}64{]} & {[}S{]} \\
Three energy regimes: linear in operations, linear in time, entropy
production & The first two are approximations, because static power
scales with time; only the third is in units where Landauer appears
directly & {[}K{]} \\
Qualcomm acquired Modular in July 2026 & Announced 24 June, closed 29
July, about \$3.1 to 3.9 billion in stock {[}65{]} & {[}S{]} \\
MLIR has no upstream path to NVIDIA or AMD GPUs & A minimal upstream
path exists (the gpu, nvvm and rocdl dialects); production-grade paths
are downstream {[}66{]} & {[}S{]}, {[}K{]} \\
NVIDIA acquired Groq for about \$20 billion (December 2025) & A
non-exclusive licence plus hires; Groq continues independently {[}67{]}
& {[}S{]} \\
TensorRT-LLM 1.2 (March 2026) removed the TensorRT backend and made
PyTorch the default & PyTorch is the sole backend; the March date was
not found {[}26{]} & {[}S{]}, {[}NOT FOUND{]} \\
CuTe-based FlashAttention-4 is 1.6× to 3.2× faster on Blackwell & That
range is FlexAttention's FlashAttention-4 backend against its own Triton
backend (GB200, forward pass); FlashAttention-4 itself is up to 1.3×
over cuDNN and 2.7× over Triton on B200 {[}68{]} & {[}S{]} \\
\end{longtable}
}

\subsection{A.3 Confirmed claims}\label{a.3-confirmed-claims}

Twenty claims were confirmed, nearly all against snippet text {[}S{]}:
``Compiling by Proving'' (September 2025) {[}24{]}; CktFormalizer's
``the compiler is the verifier'', from a third-party review only
{[}45{]}; Veritas, correct only relative to its CNF clauses {[}44{]};
EMI testing finding more than 1,000 bugs {[}2{]}; the MLGO gains of
0.3\% to 1.5\% in queries per second and 3\% to 7\% in size {[}37{]};
SuperCoder's 95\% and 1.46×, judged by tests, not proof {[}41{]}; the
Landauer figure of about 2.87 zeptojoules at 300 K {[}K, arithmetic{]};
the PP-completeness abstract, with the caveat in §4.2 {[}4{]}; the June
2026 codon-optimisation paper on thermodynamic hardware {[}69{]}; the
4\% to 886\% energy cost of model self-refinement {[}33{]}; TypeScript
7's 8× to 12× faster builds, from Microsoft's own benchmarks {[}70{]};
Mojo 1.0 on 11 August 2026 {[}71{]}; WebAssembly 3.0 in September 2025
{[}72{]}; the OctoAI, CentML, Brium and Fermyon acquisitions {[}73{]};
ORNL's JACC, from ORNL's own sources {[}74{]}; LFortran compiling LAPACK
and fpm, while still alpha {[}75{]}; gcobol in GCC 15 and 16 {[}76{]};
AdaCore's GNAT Foundry: Intersection, from secondary sources {[}39{]};
FlexAttention at about 60\% of FlashAttention-3 on Hopper {[}77{]}; and
Triton's BF16 GEMM at 62\% of cuBLAS on B200, the low end of 62\% to
70\% {[}78{]}.

\subsection{A.4 Claims not found}\label{a.4-claims-not-found}

No source was located for the following seventeen claims. None of them
is stated as fact anywhere in this study.

{\footnotesize\setlength{\tabcolsep}{3.8pt}\renewcommand{\arraystretch}{1.15}
\begin{longtable}[]{@{}
  >{\raggedright\arraybackslash}p{(\linewidth - 2\tabcolsep) * \real{0.5104}}
  >{\raggedright\arraybackslash}p{(\linewidth - 2\tabcolsep) * \real{0.4896}}@{}}
\toprule\noalign{}
\begin{minipage}[b]{\linewidth}\raggedright
\textbf{Claim}
\end{minipage} & \begin{minipage}[b]{\linewidth}\raggedright
\textbf{What was found instead}
\end{minipage} \\
\midrule\noalign{}
\endhead
\bottomrule\noalign{}
\endlastfoot
Debian is over 95\% reproducible, and its ``forky'' release makes
reproducibility a hard requirement & One machine at 97.38\% {[}79{]} \\
AutoPhase is 16\% better than -O3 & AutoPhase targets cycle counts in
high-level synthesis; its figures were not checked \\
``AI as a Compiler'' (September 2026), 0.8× to 3.34× over autotuned
Triton & A summary of a Triton-to-PTX system with a verifier reports
1.37× and 1.10× {[}80{]} \\
An October 2026 hybrid-automation study cut model-usage cost by 48\% &
Nearby figures measure other things \\
Synapze: \$2.03 against \$8 per correct application & Nothing
relevant \\
G1 uses Z3-certified rewrites & Not matched to any system; Alive2 and
Souper check LLVM rewrites with SMT, while egg does not certify its
rules \\
HPE Labs (January 2026), ``Thermodynamically Optimal Compilation with
Explicit Convergence Rates'', 99.93\% & Only unreviewed, self-published
preprints {[}36{]} \\
Extropic's 146× and a product named ``Z1/T'' & The public figure is
about 10,000×, from a GPU simulation {[}81{]} \\
Zig's self-hosted backend: 10 s and 2.8 GB against 75 s and 9.6 GB on
LLVM & One user reports 15 s against 84 s; the direction holds \\
67\% of respondents use Wasm in production & 41\%, from a 2023 survey
{[}82{]} \\
Julia 1.14 cuts load and precompile time by 30\% to 50\% & The 30\%
figure belongs to 1.13 {[}83{]} \\
LLVM Flang is production-ready on NERSC workloads & NERSC warns that
LLVM ``isn't fully compatible with all vendor software'' {[}84{]} \\
GnuCOBOL 4.0 & Version 3.2 (July 2023) is current {[}85{]} \\
TIRx is within 4\% of cuBLAS on FP16 matrix multiplication & The TIRx
post exists; no benchmark figure was visible {[}86{]} \\
A September 2026 survey: hardware outruns portable abstractions, with
three workarounds & Nearest is an August 2026 essay, ``Retire the
Abstractions'' {[}87{]} \\
MLIR has 500 to 800 commits a month & A proxy count gives 225 to 458 a
month; the exact count was refused {[}58{]} \\
A Landauer floor and a Kolmogorov target are parts of an established
compiler law & No compiler-correctness or compiler-design source ties
either to a compiler law \\
\end{longtable}
}

\section{References}\label{references}

Each entry gives the evidence mark for how the source was checked.
{[}S{]} entries were seen only as search results or secondary summaries;
the paper itself was not read.

\begin{enumerate}
\def\labelenumi{\arabic{enumi}.}
\tightlist
\item
  {[}1{]} X. Yang, Y. Chen, E. Eide and J. Regehr, ``Finding and
  Understanding Bugs in C Compilers,'' PLDI 2011.
  \href{https://www-old.cs.utah.edu/~regehr/papers/pldi11-preprint.pdf}{Preprint}
  Seen only as a search-result snippet or secondary summary {[}S{]}.
\item
  {[}2{]} C. Sun, V. Le, Z. Su and colleagues, EMI project page
  (equivalence modulo inputs), University of California, Davis.
  \url{https://web.cs.ucdavis.edu/~su/emi-project/} Seen only as a
  search-result snippet or secondary summary {[}S{]}.
\item
  {[}3{]} Zig Software Foundation, Zig 0.15.1 release notes.
  \url{https://ziglang.org/download/0.15.1/release-notes.html} Seen only
  as a search-result snippet or secondary summary {[}S{]}. The host
  could not be reached from the research sessions.
\item
  {[}4{]} A. Condon, N. Kornerup and co-authors, ``The Thermodynamic
  Complexity of Circuits,'' arXiv:2610.33193, October 2026.
  \url{https://arxiv.org/abs/2610.33193} Seen only as a search-result
  snippet or secondary summary {[}S{]}. Only the abstract was seen;
  whether the problem is in PP must be checked in the full text.
\item
  {[}5{]} USLegal, ``Compilation'' (legal definition).
  \url{https://definitions.uslegal.com/c/compilation/} Seen only as a
  search-result snippet or secondary summary {[}S{]}.
\item
  {[}6{]} openIE-dev/rce, \texttt{compile/}: CHARTER.md, HYPOTHESES.md
  and the experiments.
  \url{https://github.com/openIE-dev/rce/tree/main/compile} The
  project's own repository.
\item
  {[}7{]} liora.io, ``PyTorch's FlexAttention with FlashAttention-4.''
  \url{https://liora.io/en/pytorchs-flexattention-with-flashattention-4-is-a-game-changer}
  Seen only as a search-result snippet or secondary summary {[}S{]}.
  Secondary source.
\item
  {[}8{]} ``Taming Bitwise Behavior in GPU Kernels with Tensor Core:
  Black-Box Reconstruction, Compiler Enforcement, and Static
  Verification,'' arXiv:2609.11356, September 2026.
  \url{https://arxiv.org/pdf/2609.11356} Seen only as a search-result
  snippet or secondary summary {[}S{]}. Provisional: only a partly
  garbled conclusion was visible.
\item
  {[}9{]} ``MMA-Sim,'' a bit-accurate model of tensor and matrix cores,
  arXiv:2511.10909. \url{https://arxiv.org/html/2511.10909v1} Seen only
  as a search-result snippet or secondary summary {[}S{]}.
\item
  {[}10{]} D. Patterson and A. Ahmed, ``The Next 700 Compiler
  Correctness Theorems (Functional Pearl),'' ICFP 2019.
  \href{https://icfp19.sigplan.org/details/icfp-2019-papers/35/The-Next-700-Compiler-Correctness-Theorems-A-Functional-Pearl-}{ICFP
  2019 page} Seen only as a search-result snippet or secondary summary
  {[}S{]}.
\item
  {[}11{]} M. Patrignani, A. Ahmed and D. Clarke, ``Formal Approaches to
  Secure Compilation: A Survey of Fully Abstract Compilation and Related
  Work,'' ACM Computing Surveys 51(6), 2019.
  \href{https://dblp.org/rec/journals/csur/PatrignaniAC19.html}{DBLP
  record} Seen only as a search-result snippet or secondary summary
  {[}S{]}.
\item
  {[}12{]} Rust project, rustc dev guide, ``Overview of the compiler.''
  \url{https://rustc-dev-guide.rust-lang.org/overview.html} Cited from
  memory; the page was not opened in the research sessions {[}K{]}.
\item
  {[}13{]} Go project, cmd/compile README.
  \url{https://go.dev/src/cmd/compile/README} Cited from memory; the
  page was not opened in the research sessions {[}K{]}.
\item
  {[}14{]} L. Bryant (supervised by W. J. Bowman), ``Compilation as
  Normalization: A Multi-language Semantics Approach to Compiler
  Correctness,'' MSc thesis, University of British Columbia, November
  2024.
  \href{https://open.library.ubc.ca/collections/24/items/1.0445241}{UBC
  Open Collections} Seen only as a search-result snippet or secondary
  summary {[}S{]}. The host could not be reached; the abstract was not
  read verbatim.
\item
  {[}15{]} N. P. Lopes, J. Lee, C.-K. Hur, Z. Liu and J. Regehr,
  ``Alive2: Bounded Translation Validation for LLVM,'' PLDI 2021.
  \href{https://www-old.cs.utah.edu/~regehr/alive2-pldi21.pdf}{PDF} Seen
  only as a search-result snippet or secondary summary {[}S{]}.
\item
  {[}16{]} Y. K. Tan, M. O. Myreen, R. Kumar, A. Fox, S. Owens and M.
  Norrish, ``The Verified CakeML Compiler Backend,'' Journal of
  Functional Programming 29, 2019.
  \href{https://kar.kent.ac.uk/71304/}{Kent repository} Seen only as a
  search-result snippet or secondary summary {[}S{]}.
\item
  {[}17{]} Bytecode Alliance, ``Cranelift: an acyclic e-graph mid-end''
  (accepted RFC).
  \url{https://github.com/bytecodealliance/rfcs/blob/main/accepted/cranelift-egraph.md}
  Cited from memory; the page was not opened in the research sessions
  {[}K{]}.
\item
  {[}18{]} Y. Yang, P. M. Phothilimthana, Y. R. Wang, M. Willsey, S. Roy
  and J. Pienaar, ``Equality Saturation for Tensor Graph
  Superoptimization,'' MLSys 2021.
  \href{https://research.google/pubs/equality-saturation-for-tensor-graph-superoptimization/}{Google
  Research page} Seen only as a search-result snippet or secondary
  summary {[}S{]}.
\item
  {[}19{]} ``Mind the Abstraction Gap: Bringing Equality Saturation to
  Real-World ML Compilers,'' OOPSLA 2025.
  \href{https://2025.splashcon.org/details/OOPSLA/42/Mind-the-Abstraction-Gap-Bringing-Equality-Saturation-to-Real-World-ML-Compilers}{SPLASH
  2025 page} Seen only as a search-result snippet or secondary summary
  {[}S{]}.
\item
  {[}20{]} Formal performance and compile-time guarantees for an
  inlining heuristic, mechanised in Rocq, arXiv:2608.20137, August 2026.
  \url{https://arxiv.org/pdf/2608.20137} Seen only as a search-result
  snippet or secondary summary {[}S{]}.
\item
  {[}21{]} Rust project, ``Enabling the next-generation trait solver on
  nightly,'' Rust blog, 21 August 2026.
  \url{https://blog.rust-lang.org/2026/08/21/enabling-next-solver-on-nightly}
  Seen only as a search-result snippet or secondary summary {[}S{]}.
\item
  {[}22{]} AdaCore, ``Migrating C to Ada/SPARK with agentic AI: a
  practical playbook,'' March 2026.
  \url{https://www.adacore.com/blog/migrating-c-to-ada-spark-with-agentic-ai-a-practical-playbook}
  Seen only as a search-result snippet or secondary summary {[}S{]}.
\item
  {[}23{]} Go project, cmd/go documentation, ``Build and test caching.''
  \url{https://pkg.go.dev/cmd/go\#hdr-Build_and_test_caching} Cited from
  memory; the page was not opened in the research sessions {[}K{]}.
\item
  {[}24{]} Zhao, Hildenbrandt, Conejero and Zhao, ``Compiling by
  Proving: Language-Agnostic Automatic Optimization from Formal
  Semantics,'' arXiv:2509.21793, September 2025.
  \url{https://arxiv.org/abs/2509.21793v1} Seen only as a search-result
  snippet or secondary summary {[}S{]}.
\item
  {[}25{]} D. A. Wheeler, ``Fully Countering Trusting Trust through
  Diverse Double-Compiling,'' PhD dissertation, George Mason University,
  2009. \href{https://arxiv.org/abs/1004.5534}{arXiv:1004.5534} Seen
  only as a search-result snippet or secondary summary {[}S{]}.
\item
  {[}26{]} NVIDIA, TensorRT-LLM, ``Migration Guide: TensorRT Backend
  Removed.''
  \url{https://nvidia.github.io/TensorRT-LLM/legacy/tensorrt-backend-removal.html}
  Seen only as a search-result snippet or secondary summary {[}S{]}.
\item
  {[}27{]} TOP500, Green500 list, June 2026.
  \url{https://top500.org/lists/green500/green500-june-2026/} Seen only
  as a search-result snippet or secondary summary {[}S{]}.
\item
  {[}28{]} J. Pallister, S. J. Hollis and J. Bennett, ``Identifying
  Compiler Options to Minimise Energy Consumption for Embedded
  Platforms,'' The Computer Journal, 2015.
  \href{https://arxiv.org/pdf/1303.6485}{arXiv:1303.6485} Seen only as a
  search-result snippet or secondary summary {[}S{]}.
\item
  {[}29{]} Comparison of gfortran, ifx and ifort energy use, Cloud
  Computing and Data Science (Wiser Pub), 2025.
  \url{https://ojs.wiserpub.com/index.php/CCDS/article/view/6280} Seen
  only as a search-result snippet or secondary summary {[}S{]}.
\item
  {[}30{]} J. S. Tyagi, J. Lynch and E. D. Demaine, ``Toward an Energy
  Efficient Language and Compiler for (Partially) Reversible
  Algorithms,'' Reversible Computation 2016.
  \href{https://arxiv.org/pdf/1605.08475}{arXiv:1605.08475} Seen only as
  a search-result snippet or secondary summary {[}S{]}.
\item
  {[}31{]} J. Hong, B. Lambson, S. Dhuey and J. Bokor, ``Experimental
  test of Landauer's principle in single-bit operations on nanomagnetic
  memory bits,'' Science Advances 2, 2016; the article also reviews the
  colloidal-bead test of Bérut et al., Nature 483, 2012.
  \href{https://pmc.ncbi.nlm.nih.gov/articles/PMC4795654}{PMC4795654}
  Seen only as a search-result snippet or secondary summary {[}S{]}.
\item
  {[}32{]} M. Horowitz, ``Computing's Energy Problem (and what we can do
  about it),'' ISSCC 2014.
  \href{https://pages.cs.wisc.edu/~markhill/restricted/isscc2014_horowitz_power_scaling.pdf}{Slides}
  Seen only as a search-result snippet or secondary summary {[}S{]}. The
  ratios to the Landauer bound are arithmetic done in the research
  sessions {[}K{]}.
\item
  {[}33{]} ``Green Code or Costly Code? Evaluating the Energy Costs of
  Iterative LLM-Based Code Optimizations,'' arXiv:2610.32891, October
  2026. \url{https://arxiv.org/abs/2610.32891} Seen only as a
  search-result snippet or secondary summary {[}S{]}.
\item
  {[}34{]} D. H. Wolpert and A. Kolchinsky, ``The thermodynamics of
  computing with circuits,'' arXiv:1806.04103.
  \url{https://arxiv.org/abs/1806.04103v2} Seen only as a search-result
  snippet or secondary summary {[}S{]}.
\item
  {[}35{]} ``Minimal thermodynamic cost of computing with circuits,''
  arXiv:2504.04031. \url{https://arxiv.org/html/2504.04031v3} Seen only
  as a search-result snippet or secondary summary {[}S{]}.
\item
  {[}36{]} ``On the Rate of Convergence to the Landauer Limit,''
  self-published Zenodo record, March 2026.
  \href{https://zenodo.org/records/18922172}{Zenodo 18922172} Seen only
  as a search-result snippet or secondary summary {[}S{]}. Unreviewed.
  Listed so that it is not mistaken for support.
\item
  {[}37{]} Google Research, ``MLGO: a machine learning framework for
  compiler optimization,'' 2022.
  \href{https://www.research.google/blog/mlgo-a-machine-learning-framework-for-compiler-optimization/}{Google
  Research blog} Seen only as a search-result snippet or secondary
  summary {[}S{]}.
\item
  {[}38{]} AlphaSignal, report on model-written Rocq proofs of
  crash-freedom for 239 OSS-Fuzz entry points, October 2026.
  \url{https://alphasignal.ai/news/anthropic-s-claude-proves-239-open-source-libraries-crash-free-exposing-fuzzing}
  Seen only as a search-result snippet or secondary summary {[}S{]}.
  Secondary source; no primary post was found.
\item
  {[}39{]} Wevolver, ``Guarding AI with deterministic trust'' (AdaCore
  GNAT Foundry: Intersection), 2026.
  \url{https://www.wevolver.com/article/guarding-ai-with-deterministic-trust}
  Seen only as a search-result snippet or secondary summary {[}S{]}.
  Secondary source.
\item
  {[}40{]} T. Tunguz, ``Hybrid state-machine agents.''
  \url{https://tomtunguz.com/hybrid-state-machine-agents/index.md} Seen
  only as a search-result snippet or secondary summary {[}S{]}. A
  classification, not a measurement.
\item
  {[}41{]} A. Wei et al., ``SuperCoder,'' arXiv:2505.11480, v3 January
  2026. \url{https://arxiv.org/pdf/2505.11480} Seen only as a
  search-result snippet or secondary summary {[}S{]}.
\item
  {[}42{]} Trooskens, Karlsberg et al., ``Compiled AI: Deterministic
  Code Generation for LLM-Based Workflow Automation,'' arXiv:2604.05150,
  April 2026. \url{https://arxiv.org/abs/2604.05150} Seen only as a
  search-result snippet or secondary summary {[}S{]}.
\item
  {[}43{]} ``Progressive Crystallization,'' arXiv:2607.07052, July 2026.
  \url{https://arxiv.org/pdf/2607.07052} Seen only as a search-result
  snippet or secondary summary {[}S{]}.
\item
  {[}44{]} ``Veritas: Deterministic Verilog Code Synthesis from
  LLM-Generated Conjunctive Normal Form,'' arXiv:2506.00005, 2025.
  \url{https://arxiv.org/html/2506.00005} Seen only as a search-result
  snippet or secondary summary {[}S{]}.
\item
  {[}45{]} Moonlight, review of ``CktFormalizer: Autoformalization of
  Natural Language into Circuit Representations.''
  \url{https://www.themoonlight.io/de/review/cktformalizer-autoformalization-of-natural-language-into-circuit-representations}
  Seen only as a search-result snippet or secondary summary {[}S{]}.
  Third-party review only; no arXiv record or author list was found.
\item
  {[}46{]} F. L. Morris, ``Advice on Structuring Compilers and Proving
  Them Correct,'' POPL 1973, pp.~144--152.
  \href{https://bibtex.github.io/POPL-1973-Morris73a.html}{Bibliographic
  record} Seen only as a search-result snippet or secondary summary
  {[}S{]}. The paper itself was not read.
\item
  {[}47{]} J. W. Thatcher, E. G. Wagner and J. B. Wright, ``More on
  Advice on Structuring Compilers and Proving Them Correct,''
  Theoretical Computer Science, 1981.
  \href{https://link.springer.com/content/pdf/10.1007/3-540-10250-7_22.pdf}{Springer}
  Seen only as a search-result snippet or secondary summary {[}S{]}.
\item
  {[}48{]} Y. Futamura, ``Partial Evaluation of Computation Process: An
  Approach to a Compiler-Compiler,'' 1971; English reprint in
  Higher-Order and Symbolic Computation 12, 1999. Summarised in
  \href{https://en.wikipedia.org/wiki/Partial_evaluation}{Wikipedia,
  ``Partial evaluation''} Seen only as a search-result snippet or
  secondary summary {[}S{]}. The original was not read; sources disagree
  on its 1971 venue.
\item
  {[}49{]} A. P. Ershov, ``Mixed computation: potential applications and
  problems for study,'' Theoretical Computer Science 18, 1982,
  pp.~41--67.
  \href{https://ftp.math.utah.edu/pub/tex/bib/idx/tcs1980/18/1/41_67.html}{Bibliographic
  record} Seen only as a search-result snippet or secondary summary
  {[}S{]}.
\item
  {[}50{]} C. A. R. Hoare, J. He and A. Sampaio, ``Normal form approach
  to compiler design,'' Acta Informatica 30, 1993. Not read {[}K{]}.
  Must be read before any claim of priority.
\item
  {[}51{]} P. Cousot and R. Cousot, ``Abstract Interpretation: A Unified
  Lattice Model for Static Analysis of Programs by Construction or
  Approximation of Fixpoints,'' POPL 1977.
  \href{https://cs.nyu.edu/~pmc309/COUSOTpapers/POPL77.shtml}{Authors'
  page} Seen only as a search-result snippet or secondary summary
  {[}S{]}.
\item
  {[}52{]} C. Lattner et al., ``MLIR: A Compiler Infrastructure for the
  End of Moore's Law,'' arXiv:2002.11054, 2020.
  \url{https://arxiv.org/pdf/2002.11054} Seen only as a search-result
  snippet or secondary summary {[}S{]}.
\item
  {[}53{]} M. Willsey, C. Nandi, Y. R. Wang, O. Flatt, Z. Tatlock and P.
  Panchekha, ``egg: Fast and Extensible Equality Saturation,'' POPL
  2021. \href{https://par.nsf.gov/servlets/purl/10224915}{NSF PAR copy}
  Seen only as a search-result snippet or secondary summary {[}S{]}.
\item
  {[}54{]} A. P. Ershov, ``On the essence of compilation,'' IFIP Working
  Conference on Formal Description of Programming Concepts, 1977. Not
  read {[}K{]}. Must be read before any claim of priority.
\item
  {[}55{]} R. Cox, ``Perfectly Reproducible, Verified Go Toolchains,''
  Go blog, 28 August 2023. \url{https://go.dev/blog/rebuild} Seen only
  as a search-result snippet or secondary summary {[}S{]}.
\item
  {[}56{]} ``SlimeCompiler: Commutative Ring Analysis for Compiler
  Optimization and Quality Assurance,'' self-published Zenodo record,
  December 2025. \href{https://zenodo.org/records/17946400}{Zenodo
  17946400} Seen only as a search-result snippet or secondary summary
  {[}S{]}. Unreviewed.
\item
  {[}57{]} Dataconomy, ``Qualcomm open-sources Modular AI platform,'' 19
  August 2026.
  \url{https://dataconomy.com/2026/08/19/qualcomm-open-sources-modular-ai-platform-challenge-nvidia/}
  Seen only as a search-result snippet or secondary summary {[}S{]}.
\item
  {[}58{]} GitHub search API, commit counts for triton-lang/triton and
  llvm/llvm-project by month, April to October 2026, queried with
  \texttt{gh\ api\ search/\allowbreak{}commits} on 1 November 2026. Read in the
  source {[}F{]}. The index covers the default branch and can lag, so
  the counts are lower bounds.
\item
  {[}59{]} POPL 2026, event listing (PriSC 2026 talk ``WP-Preserving
  Compilation,'' Abate, Elsheikh et al.).
  \url{https://popl26.sigplan.org/search/all/events} Seen only as a
  search-result snippet or secondary summary {[}S{]}.
\item
  {[}60{]} Liang, Stone, Tian et al., ``Learning Compiler Pass Orders
  using Coreset and Normalized Value Prediction,'' ICML 2023.
  \href{https://arxiv.org/abs/2301.05104}{arXiv:2301.05104} Seen only as
  a search-result snippet or secondary summary {[}S{]}.
\item
  {[}61{]} AI Weekly, ``Anthropic and Theorem use Claude to verify 123
  OSS libraries,'' October 2026.
  \url{https://aiweekly.co/alerts/anthropic-theorem-use-claude-to-verify-123-oss-libraries}
  Seen only as a search-result snippet or secondary summary {[}S{]}.
\item
  {[}62{]} Anthropic, research index.
  \url{https://www.anthropic.com/research} Read in the source {[}F{]}.
  It lists no October 2026 post on this work.
\item
  {[}63{]} S. Whitelam, ``Generative thermodynamic computing,'' Physical
  Review Letters 136, 037101, 2026.
  \href{https://arxiv.org/html/2506.15121v3}{arXiv:2506.15121} Seen only
  as a search-result snippet or secondary summary {[}S{]}.
\item
  {[}64{]} Normal Computing, ``CN101: A Digital Thermodynamic Computer
  for Generative AI,'' arXiv:2608.00754.
  \url{https://arxiv.org/pdf/2608.00754} Seen only as a search-result
  snippet or secondary summary {[}S{]}.
\item
  {[}65{]} RuntimeWire, report on the Qualcomm acquisition of Modular
  and the Mojo source release, 2026.
  \url{https://runtimewire.com/article/chris-lattner-open-sources-mojo-qualcomm-modular}
  Seen only as a search-result snippet or secondary summary {[}S{]}.
\item
  {[}66{]} LLVM project, MLIR ``gpu'' dialect documentation.
  \url{https://mlir.llvm.org/docs/Dialects/GPU/} Seen only as a
  search-result snippet or secondary summary {[}S{]}.
\item
  {[}67{]} AIN, ``Nvidia is acquiring the assets of startup Groq,'' 26
  December 2025.
  \url{https://en.ain.ua/2025/12/26/nvidia-is-acquiring-the-assets-of-startup-groq/}
  Seen only as a search-result snippet or secondary summary {[}S{]}.
\item
  {[}68{]} Lambda, ``FlashAttention-4 gives the NVIDIA Blackwell
  platform its most optimized attention kernel yet.''
  \url{https://lambda.ai/blog/flashattention-4-gives-the-nvidia-blackwell-platform-its-most-optimized-attention-kernel-yet}
  Seen only as a search-result snippet or secondary summary {[}S{]}.
\item
  {[}69{]} Jelincic and Walker, codon optimisation on thermodynamic
  hardware, arXiv:2606.17327, June 2026.
  \url{https://arxiv.org/abs/2606.17327} Seen only as a search-result
  snippet or secondary summary {[}S{]}.
\item
  {[}70{]} Microsoft, ``Announcing TypeScript 7.0,'' TypeScript blog, 8
  July 2026.
  \url{https://devblogs.microsoft.com/typescript/announcing-typescript-7.0/}
  Seen only as a search-result snippet or secondary summary {[}S{]}.
\item
  {[}71{]} Techzine, ``Mojo programming language reaches version 1.0,''
  August 2026.
  \url{https://www.techzine.eu/news/devops/143570/mojo-programming-language-reaches-version-1-0/}
  Seen only as a search-result snippet or secondary summary {[}S{]}.
\item
  {[}72{]} x-cmd blog, WebAssembly 3.0 announcement summary, September
  2025. \url{https://www.x-cmd.com/blog/250924} Seen only as a
  search-result snippet or secondary summary {[}S{]}.
\item
  {[}73{]} GeekWire, ``Chip giant Nvidia acquires OctoAI,'' 2024.
  \url{https://www.geekwire.com/2024/chip-giant-nvidia-acquires-octoai-a-seattle-startup-that-helps-companies-run-ai-models/}
  Seen only as a search-result snippet or secondary summary {[}S{]}.
\item
  {[}74{]} Oak Ridge National Laboratory, ``ORNL debuts JACC for
  performance-portable Julia on HPC systems,'' June 2026.
  \url{https://www.ornl.gov/news/ornl-debuts-jacc-performance-portable-julia-hpc-systems}
  Seen only as a search-result snippet or secondary summary {[}S{]}.
\item
  {[}75{]} LFortran, ``LFortran compiles LAPACK,'' February 2026.
  \url{https://lfortran.org/blog/2026/02/lfortran-compiles-lapack/} Seen
  only as a search-result snippet or secondary summary {[}S{]}.
\item
  {[}76{]} Phoronix, ``GCC 15 merges COBOL,'' 2025; Arch Linux package
  gcc-gcobol 16.1.
  \href{https://www.phoronix.com/news/GCC-15-Merges-COBOL}{Phoronix};
  \href{https://archlinux.org/packages/core/x86_64/gcc-gcobol/}{Arch}
  Seen only as a search-result snippet or secondary summary {[}S{]}.
\item
  {[}77{]} PyTorch, ``FlexAttention + FlashAttention-4: fast and
  flexible.''
  \url{https://pytorch.org/blog/flexattention-flashattention-4-fast-and-flexible/}
  Seen only as a search-result snippet or secondary summary {[}S{]}.
\item
  {[}78{]} ``Evaluating CUDA Tile for AI Workloads on Hopper and
  Blackwell GPUs,'' arXiv:2604.23466.
  \url{https://arxiv.org/pdf/2604.23466} Seen only as a search-result
  snippet or secondary summary {[}S{]}.
\item
  {[}79{]} debian-vote mailing list, May 2025 (one machine at 97.38\%
  reproducible).
  \url{https://lists.debian.org/debian-vote/2025/05/msg00123.html} Seen
  only as a search-result snippet or secondary summary {[}S{]}.
\item
  {[}80{]} Digg summary of an ``AI-as-compiler'' system translating
  Triton to PTX with a verifier. \url{https://digg.com/ai/g2dvabpv} Seen
  only as a search-result snippet or secondary summary {[}S{]}.
  Secondary source; the original could not be reached.
\item
  {[}81{]} Extropic, denoising thermodynamic models, arXiv:2510.23972.
  \url{https://arxiv.org/html/2510.23972v1} Seen only as a search-result
  snippet or secondary summary {[}S{]}.
\item
  {[}82{]} The New Stack, WebAssembly adoption survey figures (2023).
  \url{https://thenewstack.io/whos-leading-webassembly-adoption-so-far-vendors/}
  Seen only as a search-result snippet or secondary summary {[}S{]}.
\item
  {[}83{]} AlternativeTo, ``Julia 1.13 released,'' September 2026.
  \url{https://alternativeto.net/news/2026/9/julia-1-13-released-with-major-speed-gains-in-precompilation-startup-and-garbage-collection/}
  Seen only as a search-result snippet or secondary summary {[}S{]}.
\item
  {[}84{]} NERSC documentation, base compilers.
  \url{https://docs.nersc.gov/development/compilers/base/} Seen only as
  a search-result snippet or secondary summary {[}S{]}.
\item
  {[}85{]} GnuCOBOL project page. \url{https://gnucobol.sourceforge.io/}
  Seen only as a search-result snippet or secondary summary {[}S{]}.
\item
  {[}86{]} Apache TVM, ``TIRx,'' TVM blog, 22 June 2026.
  \url{https://tvm.apache.org/2026/06/22/tirx} Seen only as a
  search-result snippet or secondary summary {[}S{]}. The page could not
  be reached; no benchmark figure was visible.
\item
  {[}87{]} Hazy Research, ``Retire the Abstractions,'' 5 August 2026.
  \url{https://hazyresearch.stanford.edu/blog/2026-08-05-retire-the-abstractions}
  Seen only as a search-result snippet or secondary summary {[}S{]}.
\end{enumerate}


\end{document}
